\documentclass[10pt]{article}
\usepackage[a4paper, margin=2cm]{geometry}

\usepackage[T1]{fontenc}
\usepackage{newtxtext}
\usepackage[utf8]{inputenc}
\usepackage[british]{babel}

\usepackage{named}

\usepackage{graphicx}
\usepackage{amsfonts}
\usepackage{amssymb}
\usepackage{amsmath}
\usepackage{amsthm}
\usepackage{algorithm2e}
\usepackage{url}
\usepackage{xcolor}

\newtheorem{definition}{Definition}
\newtheorem{theorem}{Theorem}
\newtheorem{lemma}{Lemma}
\newtheorem{corollary}{Corollary}

\newcommand{\ols}[1]{\mskip.5\thinmuskip\overline{\mskip-.5\thinmuskip {#1} \mskip-.5\thinmuskip}\mskip.5\thinmuskip} % overline short
\newcommand{\olsi}[1]{\,\overline{\!{#1}}} % overline short italic

\title{\vspace{-4ex}KG-ER$_0$, \textsf{CARM}, and  \textsf{gCARM}\\ - draft -}
\author{Enrico Franconi et al}
\date{\today}
\begin{document}
\maketitle

\immediate\write18{git log -1 --format=\%ad --date=short > gitdate.tex}
\input{gitdate.tex}

%\thispagestyle{empty}

%%%%%%%%%%%%%%%%%%%%%%%%%%%%%%%%%%%%%%%%%%%%%%%%%%%%%%%%%%%%%%%%%%%%%

\section{KG-ER$_0$ %{\small(KG-ER with simple keys, i.e., tree pattern keys of depth 1)}
}

%%%

\subsection*{Vocabulary}
$\mathcal{T}$ (Object type names $T$) \\
\phantom{}\hspace{1em}
partitioned into $\mathcal{E}$ (entity type names $E$) \\
\phantom{}\hspace{1em}
and $\mathcal{V}$ (value type names $V$)\\
$\mathcal{P}$ (relationship type names $P$, with associated arity $\alpha(P)\geq 2$)\\
$\mathcal{R}$ (role names $R$)

\noindent

%%%

\subsection*{Declarations}

\noindent
$\textsc{RelType}(P\ (T_1, \cdots , T_n))$ 
~~~~~{ 
\parbox[t]{\textwidth}{%
one for each $P\in\mathcal{P}$\\
$n=\alpha(P)$\\
\parbox[t]{\textwidth}{%
induces~ 
$\tau:\{(P,i)\mid P\in\mathcal P,\ 1\leq i\leq\alpha(P)\}\to\mathcal T$ ~~~with $\tau(P,i)=T_i$}}}

\noindent
$\textsc{RoleName}(P.i, R)$ 
~~~~~~~~~~~~~~~~~~~~~~~{ 
\parbox[t]{\textwidth}{%
$1\leq i\leq \alpha(P)$\\
induces a bijection from $\mathcal{R}$ to 
$\displaystyle\{(P,i)\mid P\in\mathcal P,\ 1\leq i\leq\alpha(P)\}$
}}

%%%

\subsection*{Constraints}

$\textsc{Mand}(E, P.i)$
\hfill { $\tau(P,i)=E$}\\
%
$\textsc{Key}(P\ (P.i_1, \cdots , P.i_m))$ \hfill{ $m>0$ ~;~ \(1\leq i_j\leq\alpha(P)\) ~;~ \(i_j\neq i_k\) for $j\neq k$}\\
%
$\textsc{Id}(E\ (P_1.i_1, \cdots , P_m.i_m))$ 
\hfill { $m>0$ ~;~ \(i_j\in\{1,2\}\) ~;~ $\alpha(P_j)=2$ and $\tau(P_j,\overline{i_j})=E$ ~with \(1\leq j\leq m\) ~;}\\
\phantom{}\hfill { $\overline{i}=1 \text{ if } i=2 \text{ and }\overline{i}=2 \text{ if } i=1 ~;~ P_h\neq P_k ~\text{for } h\neq k$}\\
%
$\textsc{IsA}((E_1,\cdots,E_m)\ E)$ \hfill{ $m\geq 1$}\\
$\textsc{Disjoint}((E_1,\cdots,E_m)\ E)$ \hfill{ $m\geq 2$}\\
$\textsc{Cover}((E_1,\cdots,E_m)\ E)$ \hfill{ $m\geq 2$}

%%%

\subsection*{Semantics}

\noindent$
\begin{array}{ll}
E: & \text{a unary relation over OIDs,}\\
V: & \text{a unary relation over values,}\\
P: & \text{a relation of arity \(\alpha(P)\) over OIDs or values.}
\end{array}
$\\
OIDs and values are disjoint.
%%%
\\

%%%
\noindent
$\textsc{RelType}(P\ (T_1, \cdots , T_n))$

$P.1\subseteq T_1 ~~~\cdots~~~ P.n\subseteq T_n$

\noindent
$\textsc{Mand}(E, P.i)$

$E\subseteq P.i$

\noindent
$\textsc{Key}(P\ (P.i_1, \cdots , P.i_m))$

$P.i_1 \cdots P.i_m\rightarrow P.1\cdots P.{\alpha(P)}$

%\noindent
%$\textsc{ExtKey}(E\ (P_1.i_1, \cdots , P_m.i_m))$
%
%$P'_j=\tau_{[\ols{i_j}\rightarrow 1\,,\,\ols{i_{j+1}}\rightarrow\,j+1]}P_j
%~~;~~P=P_1'\bowtie\cdots\bowtie P'_m
%~~;~~P.2 \cdots P.m\!+\!1\rightarrow P.1~~;~~E\subseteq P.1$ % with $\tau(P,1)=E$

%exists $P$ s.t. $\tau_{[\ols{i_j}\rightarrow 1\,,\,\ols{i_{j+1}}\rightarrow\,j+1]}P_j=\pi_{1,j}P
%~~;~~P.2 \cdots P.m\!+\!1\rightarrow P.1~~;~~E\subseteq P.1$ % with $\tau(P,1)=E$

\vspace{.5ex}\noindent
$\textsc{Id}(E\ (P_1.i_1, \cdots , P_m.i_m))$

$
P_\iota.2 \cdots P_\iota.(m\!+\!1)\rightarrow P_\iota.1
\qquad
P_\iota.1\rightarrow P_\iota.2\cdots P_\iota.(m\!+\!1)
\qquad
E\subseteq P_\iota.1 
$

for a fresh auxiliary derived relation \(P_\iota\) for this constraint \(\iota\)
\textit{defined} as:

\qquad$
P_\iota=P_1'\bowtie\cdots\bowtie P'_m 
\qquad
P'_j=\rho_{[\ols{i_j}\,\rightarrow\, 1\,,\,{i_j}\,\rightarrow\, j+1]}P_j
\quad\text{with }1\leq j\leq m
$

\noindent
$\textsc{IsA}((E_1,\cdots,E_m)\ E)$

$E_1\subseteq E ~~;~~\cdots ~~;~~E_m\subseteq E$

\noindent
$\textsc{Disjoint}((E_1,\cdots,E_m)\ E)$ 
 
$E_1\subseteq E ~~;~~\cdots ~~;~~E_m\subseteq E ~~;~~
E_i\cap E_j=\emptyset \text{ with } i\neq j$

\noindent
$\textsc{Cover}((E_1,\cdots,E_m)\ E)$

$E_1\subseteq E ~~;~~\cdots ~~;~~E_m\subseteq E ~~;~~
E\subseteq E_1 \cup\cdots\cup E_m$

%%%%%%%%%%%%%%%%%%%%%%%%%%%%%%%%%%%%%%%%%%%%%%%%%%%%%%%

% \newpage
%\thispagestyle{empty}

\begin{definition}[ER-Satisfiability]
A KG-ER$_0$ schema \(\mathcal S\) is \emph{ER-satisfiable} if it is
logically satisfiable, \(\mathcal S\not\models E=\emptyset\) for each
\(E\in\mathcal E\), and \(\mathcal S\not\models P=\emptyset\) for each
\(P\in\mathcal P\).
\end{definition}


\begin{definition}[Identity resolving]
A KG-ER$_0$ schema $\mathcal S$ induces a preorder \(\preceq_{\mathcal S}\)
over entity types, where \(E_i\preceq_{\mathcal S} E\) iff
\(\mathcal S\models E_i\subseteq E\).  Collapsing the equivalence classes
of the preorder, where \(E\equiv_{\mathcal S}E'\) iff
\(E\preceq_{\mathcal S}E'\) and \(E'\preceq_{\mathcal S}E\), yields a
partial order over equivalence classes of entity types.  The Hasse
diagram of $\mathcal S$, denoted \(H_{\mathcal S}\), is the directed
graph whose edge relation is the transitive reduction of this partial
order, oriented from each class to its immediate larger classes.

A schema $\mathcal S$ is identity resolving if:
\begin{itemize}
\item the schema \(\mathcal S\) is ER-satisfiable;
\item each connected component of \(H_{\mathcal S}\) has a unique
maximal equivalence class;
\item these maximal classes are pairwise disjoint, i.e.,
\(\mathcal S\models E\cap E'=\emptyset\) for any representatives \(E\)
and \(E'\) of two distinct maximal classes;
%
\item for each entity type \(E\), the schema \(\mathcal S\) entails an
identity constraint for \(E\), not necessarily declared explicitly in
\(\mathcal S\), whose identifying positions all have value type.  That
is, for some binary relations $P_1,\cdots,P_m$ it holds that
\(\mathcal S\models\textsc{Id}(E\ (P_1.i_1,\cdots,P_m.i_m))\) with
\(\tau(P_j,i_j)\in\mathcal V\) for \(1\leq j\leq m\).
%
\end{itemize}
The roots of \(H_{\mathcal S}\), namely the unique maximal equivalence
classes of its connected components, are the \emph{kind classes} of
\(\mathcal S\).  When an entity type name is needed, any fixed
representative of a kind class is called a \emph{kind}.
\end{definition}

\begin{theorem}[Decidability]
\label{thm:kger-decidability}
The entailment problem for KG-ER$_0$ schemas and the non-emptiness
problem for entity types and relationship types are decidable and
EXPTIME-complete.
\end{theorem}

The proof is given in Appendix~A.1.

\begin{definition}[Join dependency]
Let \(R(U)\) be a relation schema. A \emph{join dependency} over \(R\)
is an expression
\(
\bowtie[X_1,\ldots,X_n],
\)
where
\(
n\geq 2,\quad X_i\subseteq U,\quad\text{and}\quad
X_1\cup\cdots\cup X_n = U.
\)
A relation \(R\) satisfies this join dependency if
\(
R =
\pi_{X_1}(R) \bowtie \cdots \bowtie \pi_{X_n}(R).
\)
%Keys, functional dependencies, and multivalued dependencies are special cases of join dependencies.
\end{definition}

\begin{definition}[5$^{th}$ normal form]
A relation $R$ is in 5$^{th}$ normal form if and only if every join
dependency of $R$ is implied by the keys of $R$.
\end{definition}

\begin{theorem}[KG-ER\(_0^+\) as a conservative extension]
\label{thm:kger-definitional-extension}
Let KG-ER\(_0^+\) be the compiled extension of KG-ER\(_0\) with a fresh
class \(\mathcal Q_{\mathsf{Id}}\) of auxiliary relationship names.
Symbols in \(\mathcal Q_{\mathsf{Id}}\) may be used only for the
relationships introduced when compiling identification constraints.  

For a KG-ER\(_0\) schema \(\mathcal S\), its KG-ER\(_0^+\) expansion
\(\mathcal S^+\) is obtained as follows.  For every identification
constraint
\(
\iota=\textsc{Id}(E\ (P_1.i_1,\ldots,P_m.i_m))
\)
introduce:
\begin{itemize}
	\item a fresh auxiliary relationship name
\(P_\iota\in\mathcal Q_{\mathsf{Id}}\) with
\(\alpha(P_\iota)=m+1\),
	\item fresh role names for its positions,
	\item $\tau(P_\iota,1)=E,$
	\item $\tau(P_\iota,j+1)=\tau(P_j,i_j) ~\text{ with } 1\leq j\leq m.$
\end{itemize}

\noindent
Replace \(\iota\) with the inclusion dependencies
\[
P_j[\ols{i_j},i_j]\subseteq P_\iota[1,j+1],
\qquad
P_\iota[1,j+1]\subseteq P_j[\ols{i_j},i_j] ~\text{ with } 1\leq j\leq m,
\qquad
E\subseteq P_\iota.1,
\]
together with the keys
\[
P_\iota.2\cdots P_\iota.(m+1)\to P_\iota.1,\qquad
P_\iota.1\to P_\iota.2\cdots P_\iota.(m+1).
\]
Then \(\mathcal S^+\) is a lossless variant of \(\mathcal S\), and 
in particular a conservative extension: every model of \(\mathcal S\) has a
unique expansion to a model of \(\mathcal S^+\), every model of
\(\mathcal S^+\) reduces to a model of \(\mathcal S\) by forgetting the
auxiliary names in \(\mathcal Q_{\mathsf{Id}}\), and for every constraint
\(\varphi\) over the original KG-ER\(_0\) vocabulary,
\[
\mathcal S\models\varphi
\quad\text{iff}\quad
\mathcal S^+\models\varphi .
\]
\end{theorem}

The proof is given in Appendix~A.1.


By Theorem~\ref{thm:kger-definitional-extension}, in the sequel we may work
in the compiled KG-ER\(_0^+\) expansion of a KG-ER\(_0\) schema and
suppress the superscript \(+\) whenever no confusion arises.

\begin{definition}[Key/non-key decomposition of KG-ER\(_0^+\)]
Let \(\mathcal S\) be a compiled KG-ER\(_0^+\) schema.  Keeping its
vocabulary and declarations fixed, write
\[
\mathcal S=\mathcal B_{\mathcal S}\cup\mathcal D_{\mathcal S}\cup
\mathcal K_{\mathcal S}\cup\mathcal N_{\mathcal S},
\]
where the four components are defined as follows:
\begin{itemize}
\item \(\mathcal B_{\mathcal S}\) is the declaration background.  For
every declared relationship position \(P.i\), including auxiliary
positions \(P_\iota.i\), it fixes the arity \(\alpha(P)\), the type
assignment \(\tau(P,i)\), the typing inclusion
\(P.i\subseteq\tau(P,i)\), and the disjointness between OIDs and
values.
\item \(\mathcal D_{\mathcal S}\) contains exactly the inclusion
dependencies introduced by the KG-ER\(_0^+\) encoding of identification
constraints.  Thus, for every compiled identification constraint
\(\iota=\textsc{Id}(E\ (P_1.i_1,\ldots,P_m.i_m))\) and every
\(1\leq j\leq m\), it contains the two projection inclusions
\[
P_j[\ols{i_j},i_j]\subseteq P_\iota[1,j+1],
\qquad
P_\iota[1,j+1]\subseteq P_j[\ols{i_j},i_j],
\]
together with the unary inclusion \(E\subseteq P_\iota.1\).  It contains
no key dependency.
\item \(\mathcal K_{\mathcal S}\) contains all key dependencies,
including the explicit keys on the auxiliary relationships \(P_\iota\);
in particular, the two keys introduced for each \(P_\iota\) belong to
\(\mathcal K_{\mathcal S}\).
\item \(\mathcal N_{\mathcal S}\) contains all remaining non-key
constraints.
\end{itemize}
\end{definition}

\begin{theorem}[Relative key/non-key entailment separation]
\label{thm:kger-pure-key-separation}
Let \(\mathcal S\) be the compiled KG-ER\(_0^+\) expansion of an
ER-satisfiable KG-ER\(_0\) schema.  Then, for every key dependency
\(\kappa\) over a non-auxiliary
relationship \(R\in\mathcal P\),
\[
\mathcal B_{\mathcal S}\cup\mathcal D_{\mathcal S}\cup
\mathcal K_{\mathcal S}\cup\mathcal N_{\mathcal S}\models\kappa
\quad\text{iff}\quad
\mathcal B_{\mathcal S}\cup\mathcal D_{\mathcal S}\cup
\mathcal K_{\mathcal S}\models\kappa .
\]
Thus, once the inclusion dependencies introduced by the KG-ER\(_0^+\)
encoding of identification constraints are fixed, the remaining non-key
constraints do not contribute to the entailment of keys.
\end{theorem}

The proof is given in Appendix~A.1.


\begin{theorem}[5$^{th}$ normal form of KG-ER$_0$]
\label{thm:kger-5}
Let \(\mathcal S\) be an ER-satisfiable KG-ER\(_0\) schema, and let
\(R\) be a declared relationship type of arity \(n\).  Let
\(\Sigma_R\) be the set of all key dependencies
\(X\to R.1\cdots R.n\) entailed by \(\mathcal S\), with
\(X\subseteq\{R.1,\ldots,R.n\}\).  Then, for every join
dependency \(J\) over the attributes of \(R\),
\[
\mathcal S\models J
\quad\text{iff}\quad
\Sigma_R\models J .
\]
Hence every such relationship type \(R\) is in 5$^{th}$ normal form.
\end{theorem}

The proof is given in Appendix~A.1.


\begin{corollary}[KG-ER$_0$ cannot express symmetric binary relations]
\label{cor:kger-no-symmetry}
Let \(\mathcal S\) be an ER-satisfiable KG-ER\(_0\) schema, let
\(R\) be a declared relationship type of arity \(2\), and assume that
\(\mathcal S\) admits a model \(M\) containing an \(R\)-tuple
\((\alpha,\beta)\) with \(\alpha\neq\beta\); equivalently,
\(\mathcal S\not\models\forall x,y.\ R(x,y)\to x=y\).
Then \(\mathcal S\) does not entail that \(R\) is symmetric:
\[
\mathcal S\not\models \forall x,y.\ R(x, y) \rightarrow R(y, x).
\]
\end{corollary}

The proof is given in Appendix~A.1.



%%%%%%%%%%%%%%%%%%%%%%%%%%%%%%%%%%%%%%%%%%%%%%%%%%%%%%%

% \newpage
%\thispagestyle{empty}
\section{\textsf{CARM}}

\begin{definition}[Dependency terminology]
For a dependency of the form \(\mathbf A \to \operatorname{att}(R)\) (key), \(\mathbf A \to \mathbf B\) (functional dependency), or
\(\mathbf A \twoheadrightarrow \mathbf B\) (multivalued dependency), we call \(\mathbf A\) the
\emph{determinant}.

For an inclusion dependency
\(
R[\mathbf A] \subseteq S[\mathbf B],
\)
with $\mathbf A\subseteq\operatorname{att}(R)$ and $\mathbf B\subseteq\operatorname{att}(S)$,
we call \((R,\mathbf A)\) (the source) and \((S,\mathbf B)\) (the target)
\emph{projection occurrences}.
\end{definition}

\begin{definition}[\textsf{CARM} relational database schema]
A \textsf{CARM} (Canonical Abstract Relational Model) relational database schema is a tuple
\[
\mathcal C =
(\mathcal Q,\operatorname{att},\operatorname{dom},
 \Sigma_K,\Sigma_I,\Sigma_{\mathsf{cov}},\Sigma_{\mathsf{disj}})
\]
satisfying the following conditions.

\begin{itemize}

\item \textbf{Relations and attributes.}
For every relation symbol \(R\in\mathcal Q\),
\(\operatorname{att}(R)\) is a finite set of attributes, equipped with a
fixed total order.  The function \(\operatorname{dom}\) maps every
attribute \(A\in\operatorname{att}(R)\) to its domain.

\item \textbf{Keys.}
\(\Sigma_K\) is a finite set of key dependencies.  A key dependency on
\(R\in\mathcal Q\) is written
\[
K \to \operatorname{att}(R),
\]
where \(K\subseteq \operatorname{att}(R)\) and $K\neq \emptyset$.  We call \(K\) a key of \(R\).

%\item \textbf{Fifth normal form.}
%For every relation symbol \(R\), every join dependency over
%\(\operatorname{att}(R)\) that is part of the dependency theory of \(R\)
%is implied by the keys of \(R\).  Equivalently, each relation schema
%\(R(\operatorname{att}(R))\) is in PJ/NF, or fifth normal form, with
%respect to its declared keys.
%
\item \textbf{Inclusion dependencies.}
\(\Sigma_I\) is a finite set of typed inclusion dependencies of the form
\[
R[\mathbf A] \subseteq S[\mathbf B]\quad\text{ with }\quad
R,S\in\mathcal Q.
\]
Here \(\mathbf A\) and \(\mathbf B\) are duplicate-free lists of
attributes of \(R\) and \(S\), respectively, ordered according to the
fixed attribute order of their relation:
\[
\mathbf A\subseteq \operatorname{att}(R),\quad
\mathbf B\subseteq \operatorname{att}(S),\quad
|\mathbf A|=|\mathbf B|\neq\emptyset
\]
where set notation refers to the underlying set of attributes in the
list.  Corresponding attributes have the same domain.  The semantics is
\[
\pi_{\mathbf A}(R)\subseteq \pi_{\mathbf B}(S).
\]

\item \textbf{Projection occurrences.}
A source or target of an inclusion dependency, covering constraint, or
disjointness constraint is called a projection occurrence.  Thus an
occurrence is a pair \((R,\mathbf A)\), where \(\mathbf A\) is an ordered
duplicate-free list of attributes of \(R\).  When no confusion is
possible, we write only \(\mathbf A\) for the occurrence
\((R,\mathbf A)\), and we treat them as ordered sets of attributes.

\item \textbf{Key/occurrence compatibility.}
For every projection occurrence \((R,\mathbf A)\) and every declared key \(K\) of
\(R\), one of the following holds:
\[
\mathbf A = K,
\qquad
\mathbf A \subsetneq K,
\qquad
\mathbf A \cap K = \emptyset.
\]
That is, no projection occurrence partially overlaps a key unless it is
contained in that key.

\item \textbf{Occurrence non-overlap.}
For any two projection occurrences \((R,\mathbf A)\) and
\((R,\mathbf B)\), either
\[
\mathbf A=\mathbf B\quad \text{ or }\quad\mathbf A\cap\mathbf B=\emptyset.
\]
That is, distinct projection occurrences of the same relation are either
identical or disjoint.

\item \textbf{Covering constraints.}
\(\Sigma_{\mathsf{cov}}\) is a finite set of constraints of the form

\(
\mathsf{Cover}
\bigl((R_1,\mathbf A_1),\ldots,(R_n,\mathbf A_n);
(S,\mathbf B)\bigr)\\
\text{with}\quad
n\geq 1,\quad R_i,S\in\mathcal Q,\quad
\mathbf A_i\subseteq \operatorname{att}(R_i),\quad
\mathbf B\subseteq \operatorname{att}(S),\quad
|\mathbf A_i|=|\mathbf B|\neq\emptyset
\)

where corresponding attributes of \(\mathbf A_i\) and \(\mathbf B\) have
the same domain.  The semantics is
\[
\pi_{\mathbf A_i}(R_i)\subseteq \pi_{\mathbf B}(S)
\quad\text{for every }i,
\qquad
\pi_{\mathbf B}(S)
=
\bigcup_{i=1}^n \pi_{\mathbf A_i}(R_i).
\]

\item \textbf{Disjointness constraints.}
\(\Sigma_{\mathsf{disj}}\) is a finite set of constraints of the form

\(
\mathsf{Disjoint}
\bigl((R_1,\mathbf A_1),\ldots,(R_n,\mathbf A_n);
(S,\mathbf B)\bigr)\\
\text{with}\quad
n\geq 2,\quad R_i,S\in\mathcal Q,\quad
\mathbf A_i\subseteq \operatorname{att}(R_i),\quad
\mathbf B\subseteq \operatorname{att}(S),\quad
|\mathbf A_i|=|\mathbf B|\neq\emptyset
\)

where corresponding attributes of \(\mathbf A_i\) and \(\mathbf B\) have
the same domain.  Its semantics is that the source projections
are pairwise disjoint:
\[
\pi_{\mathbf A_i}(R_i)\subseteq \pi_{\mathbf B}(S)
\quad\text{for every }i,
\qquad
\pi_{\mathbf A_i}(R_i)\cap \pi_{\mathbf A_j}(R_j)=\emptyset
\quad\text{for } i\neq j.
\]

\item \textbf{Forest presentation.}
A schema \(\mathcal C\) induces a preorder
\(\preceq_{\mathcal C}\) over the finite set of projection
occurrences of \(\mathcal C\), where
\((R,\mathbf A)\preceq_{\mathcal C}(S,\mathbf B)\) iff
\(\mathcal C\models R[\mathbf A]\subseteq S[\mathbf B]\).  Let
\((R,\mathbf A)\equiv_{\mathcal C}(S,\mathbf B)\) iff both
\((R,\mathbf A)\preceq_{\mathcal C}(S,\mathbf B)\) and
\((S,\mathbf B)\preceq_{\mathcal C}(R,\mathbf A)\).  Quotienting by
\(\equiv_{\mathcal C}\) yields a partial order over equivalence classes
of projection occurrences.  The Hasse diagram \(H_{\mathcal C}\) is the
directed graph whose edge relation is the transitive reduction of this
partial order, oriented from each class to its immediate larger classes.

A \textsf{CARM} schema must have a unique maximal class in each connected
component of \(H_{\mathcal C}\).  We then fix a covering spanning forest
\(F_{\mathcal C}\) of \(H_{\mathcal C}\) by choosing, for every
non-maximal class, one outgoing Hasse edge to a larger class on a path to
the unique maximal class of its component.  Thus \(F_{\mathcal C}\) is
rooted towards the maximal classes; these maximal classes are exactly its
roots.
\end{itemize}

\end{definition}

\begin{theorem}[5$^{th}$ normal form of \textsf{CARM}]
\label{lem:carm-5}
Let \(\mathcal C\) be a satisfiable \textsf{CARM} schema, and let \(R\)
be a relation symbol that is non-empty in some model of \(\mathcal C\).
Let \(\Sigma_R\) be the set of all key dependencies
\(X\to\operatorname{att}(R)\) entailed by \(\mathcal C\). Then, for every
join dependency \(J\) over \(\operatorname{att}(R)\),
\[
\mathcal C\models J
\quad\text{iff}\quad
\Sigma_R\models J .
\]
Hence every such relation \(R\) is in 5$^{th}$ normal form.
\end{theorem}

The proof is given in Appendix~A.2.



%\begin{definition}[Non-conflicting keys and TGDs~\cite{DBLP:journals/ws/CaliGL12}]
%Let $\kappa$ be a key, and let $\sigma$ be a TGD of the form
%$\Phi(\mathbf{X}, \mathbf{Y}) \to \exists \mathbf{Z}\, P(\mathbf{X}, \mathbf{Z}).$
%Then $\kappa$ is \emph{non-conflicting} $(NC)$ with $\sigma$ iff either:
%\begin{enumerate}
%    \item the relational predicate on which $\kappa$ is defined is different from $P$; or
%    \item the positions of $\kappa$ in $P$ are not a proper subset of the
%    $\mathbf{X}$-positions in $P$ in the head of $\sigma$, and every variable
%    in $\mathbf{Z}$ appears only once in the head of $\sigma$.
%\end{enumerate}
%
%\noindent
%We say $\kappa$ is \emph{non-conflicting} $(NC)$ with a set of TGDs
%$\Sigma_T$ iff $\kappa$ is NC with every $\sigma \in \Sigma_T$.\\
%A set of keys $\Sigma_K$ is \emph{non-conflicting} $(NC)$ with $\Sigma_T$
%iff every $\kappa \in \Sigma_K$ is NC with $\Sigma_T$.
%\end{definition}
%
%%%

%\begin{theorem}[Decidability of query answering]
%\textsf{CARM} is a special case of a database schema with non-conflicting keys and inclusion dependencies (linear TGDs), in which the separation property between keys and inclusion dependencies holds. 
%\end{theorem}

%%%%%%%%%%%%%%%%%%%%%%%%%%%%%%%%%%%%%%%%%%%%%%%%%%%%%%%

% \newpage
\section{KG-ER$_0$ $\leftrightarrow$ \textsf{CARM} --- lossless transformations}

\begin{theorem}[KG-ER$_0$ to \textsf{CARM} with OIDs]
\label{thm:kgerzero-to-carm-oids}
Let $\mathcal S$ be a KG-ER$_0$ schema.  Then
$\mathcal S$ can be translated in linear time into a \textsf{CARM}
schema $\mathcal C_{\mathcal S}$ over the same OID and value domains.
The models of $\mathcal S$ and the models of $\mathcal C_{\mathcal S}$
are in bijection, up to the definitional expansion/reduct of the
relations introduced for \textsc{Id} constraints.
\end{theorem}

The proof is given in Appendix~A.3.


\begin{theorem}[Identity resolving KG-ER$_0$ to \textsf{CARM} without OIDs]
\label{thm:kgerzero-to-carm-no-oids}
Let $\mathcal S$ be an identity resolving KG-ER$_0$ schema.  
Then $\mathcal S$ can be transformed into an OID-free \textsf{CARM}
schema $\mathcal C^\circ_{\mathcal S}$ over value domains only.  The
models of $\mathcal S$ and the models of
$\mathcal C^\circ_{\mathcal S}$ are in bijection up to renaming of OIDs
and definitional expansion/reduct of the fixed identifier relationships.
Once the identities $\iota_K$ have been fixed, the transformation is
linear in the size of the schema.
\end{theorem}

The proof is given in Appendix~A.3.


\begin{theorem}[\textsf{CARM} without OIDs to KG-ER$_0$]\label{thm:carm-to-kgerzero}
Let $\mathcal C$ be a \textsf{CARM} schema without OID domains.  Assume
that the inclusion, covering, and disjoint constraints of $\mathcal C$
are given by their forest presentation: the nodes are ordered
source/target projection occurrences and the edges are the inclusion
dependencies between them.  Then $\mathcal C$ can be transformed into a
KG-ER$_0$ schema $\mathcal K_{\mathcal C}$ such that the instances of
$\mathcal C$ are in bijection, up to renaming of the introduced OIDs,
with the instances of $\mathcal K_{\mathcal C}$.  In particular, the
transformation is lossless.
\end{theorem}

The proof is given in Appendix~A.3.


%%%%%%%%%%%%%%%%%%%%%%%%%%%%%%%%%%%%%%%%%%%%%%%%%%%%%%%

% \newpage
%\thispagestyle{empty}
\section{\textsf{gCARM} $\leftrightarrow$ \textsf{CARM} --- lossless transformations}

\subsection*{\textsf{gCARM} relational database schema}
\begin{itemize}
\item with key, functional, and multivalued dependencies, having an extended conflict-free cover.
\item with \textit{inclusion dependencies} such that:
\begin{itemize}
	\item each source and target of each inclusion dependency is
    \begin{itemize}
	\item either equal to the source of a dependency,
	% \item or a subset of a source of a dependency, 
	\item or not overlapping with any source of a dependency,
    \end{itemize} 	
    \item and
	\begin{itemize}
	\item either equal to a different source or target of an inclusion dependency, 
	\item or not overlapping with any different source or target of an inclusion dependency.
    \end{itemize} 	
\end{itemize}
\item with \textit{covering} and \textit{disjoint} constraints over sources of inclusion dependencies with common targets.
\end{itemize}

\begin{definition}[Extended conflict-free dependencies]
Let \(U\) be a finite set of attributes, and let \(D\) be a set of
functional dependencies and multivalued dependencies over \(U\).
%
First define the envelope of \(D\) as
\[
\operatorname{Env}(D)
=
\{\, X \twoheadrightarrow Y
\mid
X \in \operatorname{LHS}(D),\;
D \models X \twoheadrightarrow Y,\;
D \not\models X \to Y
\,\}.
\]
That is, \(\operatorname{Env}(D)\) contains the genuine multivalued
dependencies implied by \(D\), excluding those that are actually functional
dependencies.
%
For a set \(M\) of multivalued dependencies, say that an MVD
\(
X \twoheadrightarrow Y
\)
splits a set of attributes \(Z \subseteq U\) if
\[
Z \cap (Y - X) \neq \varnothing
\quad\text{and}\quad
Z \cap (U - XY) \neq \varnothing.
\]

For \(X \subseteq U\), let \(\operatorname{dep}_{M}(X)\) denote the
dependency basis of \(X\) with respect to \(M\), i.e. the partition of
\(U-X\) induced by the MVDs determined by \(X\).
%
The set \(M\) is conflict-free if:
\begin{enumerate}
    \item no MVD in $M$ splits any $X \in \operatorname{LHS}(M)$,
    \item and, for all \(X,Y \in \operatorname{LHS}(M)\), 
    $\operatorname{dep}_{M}(X) \cap \operatorname{dep}_{M}(Y)
    \subseteq
    \operatorname{dep}_{M}(X \cap Y).$
\end{enumerate}
%
The set \(D\) is extended conflict-free if \(\operatorname{Env}(D)\) is conflict-free.

Finally, \(D\) has an extended conflict-free cover if there exists a set
\(D'\) of FDs and MVDs such that
$
D^+ = {D'}^+
$
and \(D'\) is extended conflict-free.
\end{definition}

\begin{theorem}[Lossless decomposition~\cite{YuanOzsoyoglu1987LogicalDesign}]
\label{thm:fdmvd}
Let \(U\) be a finite set of attributes, and let \(D\) be a set of
functional dependencies and multivalued dependencies over \(U\) --- no other join dependencies.
\(D\) admits a lossless, dependency-preserving 5NF decomposition if and only if
\(D\) has an extended conflict-free cover.
\end{theorem}

% \begin{definition}[Boyce--Codd Normal Form]
% Let \(R(U)\) be a relation schema over attributes \(U\), and let \(F\) be a
% set of functional dependencies over \(U\). We say that \(R\) is in
% Boyce--Codd normal form (BCNF) with respect to \(F\) if, for every
% functional dependency
% \(
% X \to Y
% \)
% such that
% \(
% F \models X \to Y,
% \)
% one of the following holds:
% \(
% Y \subseteq X
% \)
% or
% \(
% F \models X \to U.
% \)
% That is, every non-trivial functional dependency has a superkey as its
% left-hand side.
% \end{definition}

% \begin{definition}[Non-critical FD cover]
% Let \(U\) be a finite set of attributes, and let \(F\) be a set of
% functional dependencies over \(U\).
% %
% A canonical cover \(G\) of \(F\) is called non-critical if:

% \begin{enumerate}
%     \item \(G^+ = F^+\);
%     \item every dependency in \(G\) has the form
%     \(
%     X \to A
%     \)
%     with \(A\) a single attribute and \(A \notin X\);
%     \item \(G\) is minimal, i.e. no dependency can be removed and no
%     attribute can be removed from the left-hand side of a dependency
%     without changing the closure;
%     % \item for every dependency \(X \to A\) in \(G\), the schema
%     % \(
%     % R_{X,A} = X \cup \{A\}
%     % \)
%     % is in BCNF with respect to \(F\).
%     \item for every \(X \to A \in G\) and every non-trivial dependency
% \(
% Y \to B
% \)
% with
% \(
% Y \cup \{B\} \subseteq R_{X,A},
% \)
% if
% \(
% F \models Y \to B,
% \)
% then \(Y\) is a superkey of \(R_{X,A}\), i.e.
% \(
% F \models Y \to R_{X,A}.
% \)
% \end{enumerate}

% % Equivalently, condition (4) means that for every \(X \to A \in G\),
% % for every non-trivial dependency
% % \[
% % Y \to B
% % \]
% % such that
% % \[
% % Y \cup \{B\} \subseteq R_{X,A},
% % \]
% % if
% % \[
% % F \models Y \to B,
% % \]
% % then \(Y\) is a superkey of \(R_{X,A}\), i.e.
% % \[
% % F \models Y \to R_{X,A}.
% % \]
% \end{definition}

% \begin{theorem}[Lossless decomposition for FDs~\cite{Osborn1979CoveringBCNF}]
% Let \(U\) be a finite set of attributes, and let \(F\) be a set of
% functional dependencies over \(U\) --- no other join dependencies.
% \(F\) admits a lossless, dependency-preserving 5NF decomposition if and only if
% \(F\) has a non-critical canonical cover.
% \end{theorem}

%\begin{definition}[\textsf{CARM}-admissible decomposition]
%Let $R(U)$ be a relation of a \textsf{gCARM} schema, let $D_R$ be the
%FDs and MVDs declared on $R$, and let $\Omega_R$ be the set of ordered
%attribute lists of $R$ that occur as sources or targets of inclusion,
%covering, or disjoint constraints.  A decomposition
%\[
%  \Delta_R=\{R^1(U^1),\ldots,R^k(U^k)\}
%\]
%is \emph{\textsf{CARM}-admissible} if:
%\begin{enumerate}
%    \item it is lossless and dependency-preserving for $D_R$, and the
%    dependencies projected on each component are expressible as keys of
%    that component;
%    \item every $X\in\Omega_R$ is contained in some component $U^a$; one
%    such component, written $\operatorname{rep}_R(X)$, is fixed;
%    \item the projection-preservation inclusions of the decomposition,
%    together with the translated inclusions on the chosen representatives
%    $\operatorname{rep}_R(X)$, satisfy the \textsf{CARM} key-overlap,
%    occurrence-overlap, and lexicographic ordering conditions;
%    \item for every family of component relations
%    $r^1,\ldots,r^k$ satisfying those keys and projection-preservation
%    inclusions,
%    \[
%      \pi_{U^a}(r^1\bowtie\cdots\bowtie r^k)=r^a
%      \qquad (1\leq a\leq k).
%    \]
%\end{enumerate}
%\end{definition}

%\begin{theorem}[\textsf{gCARM} to \textsf{CARM}]
%\label{thm:gcarm-to-carm}
%Let $\mathcal G$ be a \textsf{gCARM} schema.  Assume that, for every
%relation $R$ of $\mathcal G$, a \textsf{CARM}-admissible 5$^{th}$ normal
%form decomposition $\Delta_R$ has been chosen; by
%Theorem~\ref{thm:fdmvd}, the lossless and dependency-preserving part of
%such a decomposition exists whenever the FD/MVD constraints of $R$ have
%an extended conflict-free cover.  Then $\mathcal G$ can be transformed
%into a \textsf{CARM} schema $\mathcal C_{\mathcal G}$ whose instances
%are in bijection with the instances of $\mathcal G$.  In particular, the
%transformation is lossless.
%\end{theorem}
\begin{theorem}[\textsf{gCARM} to \textsf{CARM}]
\label{thm:gcarm-to-carm}
Let $\mathcal G$ be a \textsf{gCARM} schema. By
Theorem~\ref{thm:fdmvd}, the lossless and dependency-preserving part of
such a decomposition exists whenever the FD/MVD constraints of $R$ have
an extended conflict-free cover.  Then $\mathcal G$ can be transformed
into a \textsf{CARM} schema $\mathcal C_{\mathcal G}$ whose instances
are in bijection with the instances of $\mathcal G$.
\end{theorem}

The proof is given in Appendix~A.4.


\begin{theorem}[\textsf{CARM} to \textsf{gCARM}]
\label{thm:carm-to-gcarm}
Every \textsf{CARM} schema $\mathcal C$ can be viewed as a
\textsf{gCARM} schema $\mathcal G_{\mathcal C}$ over the same relation
symbols and domains.  The two schemas have exactly the same instances;
hence the transformation is lossless.
\end{theorem}

The proof is given in Appendix~A.4.


%%%%%%%%%%%%%%%%%%%%%%%%%%%%%%%%%%%%%%%%%%%%%%%%%%%%%%%

\newpage
\section*{Appendix A.1. Proofs for Section 1}

\noindent\textbf{Statement of Theorem~\ref{thm:kger-decidability} (\colorbox{yellow}{\textbf{Decidability}}).}
The entailment problem for KG-ER$_0$ schemas and the non-emptiness
problem for entity types and relationship types are decidable and
EXPTIME-complete.

\begin{proof}[Proof of Theorem~\ref{thm:kger-decidability} (Decidability)]
EXPTIME-hardness follows by a polynomial time reduction from
\(\mathcal{ALC}\) satisfiability, and membership follows by a polynomial
time reduction to \(\mathcal{DLR}^{\pm}\)
satisfiability~\cite{DBLP:conf/birthday/ArtaleF19}.
\end{proof}

\noindent\textbf{Statement of Theorem~\ref{thm:kger-definitional-extension} (\colorbox{yellow}{KG-ER\(_0^+\) as a conservative extension}).}
Let KG-ER\(_0^+\) be the compiled extension of KG-ER\(_0\) with a fresh
class \(\mathcal Q_{\mathsf{Id}}\) of auxiliary relationship names.
Symbols in \(\mathcal Q_{\mathsf{Id}}\) may be used only for the
relationships introduced when compiling identification constraints.  

For a KG-ER\(_0\) schema \(\mathcal S\), its KG-ER\(_0^+\) expansion
\(\mathcal S^+\) is obtained as follows.  For every identification
constraint
\(
\iota=\textsc{Id}(E\ (P_1.i_1,\ldots,P_m.i_m))
\)
introduce:
\begin{itemize}
	\item a fresh auxiliary relationship name
\(P_\iota\in\mathcal Q_{\mathsf{Id}}\) with
\(\alpha(P_\iota)=m+1\),
	\item fresh role names for its positions,
	\item $\tau(P_\iota,1)=E,$
	\item $\tau(P_\iota,j+1)=\tau(P_j,i_j) ~\text{ with } 1\leq j\leq m.$
\end{itemize}

\noindent
Replace \(\iota\) with the inclusion dependencies
\[
P_j[\ols{i_j},i_j]\subseteq P_\iota[1,j+1],
\qquad
P_\iota[1,j+1]\subseteq P_j[\ols{i_j},i_j] ~\text{ with } 1\leq j\leq m,
\qquad
E\subseteq P_\iota.1,
\]
together with the keys
\[
P_\iota.2\cdots P_\iota.(m+1)\to P_\iota.1,\qquad
P_\iota.1\to P_\iota.2\cdots P_\iota.(m+1).
\]
Then \(\mathcal S^+\) is a lossless variant of \(\mathcal S\), and 
in particular a conservative extension: every model of \(\mathcal S\) has a
unique expansion to a model of \(\mathcal S^+\), every model of
\(\mathcal S^+\) reduces to a model of \(\mathcal S\) by forgetting the
auxiliary names in \(\mathcal Q_{\mathsf{Id}}\), and for every constraint
\(\varphi\) over the original KG-ER\(_0\) vocabulary,
\[
\mathcal S\models\varphi
\quad\text{iff}\quad
\mathcal S^+\models\varphi .
\]

\begin{proof}[Proof of Theorem~\ref{thm:kger-definitional-extension} (KG-ER\(_0^+\) as a conservative extension)]
Let \(M\models\mathcal S\).  Keep the interpretation of every old type
and relationship name unchanged.  For each identification constraint
\(\iota\), define
\[
P_\iota^M=(P'_1)^M\bowtie\cdots\bowtie(P'_m)^M,
\]
where
\[
P'_{j}=\rho_{[\ols{i_j}\,\rightarrow\,1,\,
                 i_j\,\rightarrow\,j+1]}P_j
\quad (1\leq j\leq m).
\]
Since \(M\) satisfies the original \(\textsc{Id}\)-clause, this relation
satisfies the two displayed keys and the inclusion
\(E\subseteq P_\iota.1\).  The inclusion
\(P_\iota[1,j+1]\subseteq P_j[\ols{i_j},i_j]\) is immediate from the
definition of the join.  Conversely, let
\((c,a)\in P_j[\ols{i_j},i_j]\).  By typing, \(c\in E\), and therefore
the inclusion \(E\subseteq P_\iota.1\) gives a tuple
\((c,b_1,\ldots,b_m)\in P_\iota\).  Since \(P_\iota\) is the natural
join \(P_1'\bowtie\cdots\bowtie P_m'\), membership of
\((c,b_1,\ldots,b_m)\) in \(P_\iota\) means, in particular, that
\((c,b_k)\in P_k[\ols{i_k},i_k]\) for every \(1\leq k\leq m\).
Together with the hypothesis \((c,a)\in P_j[\ols{i_j},i_j]\), all
\(m\) binary-projection membership conditions are met for the tuple
\((c,b_1,\ldots,b_{j-1},a,b_{j+1},\ldots,b_m)\), so this tuple also
belongs to the natural join \(P_\iota\).  The key
\(P_\iota.1\to P_\iota.2\cdots P_\iota.(m+1)\) forces \(a=b_j\).  Hence
\((c,a)\in P_\iota[1,j+1]\), proving the other projection inclusion.
Thus \(M\) expands to a model of \(\mathcal S^+\).

Conversely, let \(N\models\mathcal S^+\).  Forget the auxiliary
relationship names in \(\mathcal Q_{\mathsf{Id}}\).  All
non-identification constraints of \(\mathcal S\) are unchanged.  For each
compiled identification constraint, the inclusion dependencies give
\[
P_\iota[1,j+1]=P_j[\ols{i_j},i_j]\quad (1\leq j\leq m).
\]
The relation \(P_\iota\) is exactly the join of these binary projections.
Indeed, every tuple of \(P_\iota\) belongs to the join by the inclusions
\(P_\iota[1,j+1]\subseteq P_j[\ols{i_j},i_j]\).  Conversely, if
\((c,a_1,\ldots,a_m)\) belongs to the join, then for each \(j\) the
inclusion \(P_j[\ols{i_j},i_j]\subseteq P_\iota[1,j+1]\) gives a
\(P_\iota\)-tuple with first component \(c\) and \((j+1)\)-st component
\(a_j\).  The key
\(P_\iota.1\to P_\iota.2\cdots P_\iota.(m+1)\) forces all these tuples to
coincide, so \((c,a_1,\ldots,a_m)\in P_\iota\).  Therefore \(P_\iota\)
satisfies the original join definition; together with the two keys and
\(E\subseteq P_\iota.1\), this is exactly the original
\(\textsc{Id}\)-semantics.  Hence the reduct of \(N\) is a model of
\(\mathcal S\).

The argument also shows uniqueness of the expansion: in every model of
\(\mathcal S^+\), each auxiliary \(P_\iota\) is forced to be the join
relation used in the original semantics.  Thus expansion and reduct are
inverse on the original vocabulary.  This gives both losslessness and
preservation and reflection of entailment over the original vocabulary.
\end{proof}

\begin{lemma}[\colorbox{yellow}{Chase-level key/non-key separation for KG-ER\(_0^+\)}]
\label{lem:kger-key-nonkey-separation}
Let \(\mathcal S\) be the compiled KG-ER\(_0^+\) expansion of an
ER-satisfiable KG-ER\(_0\) schema, let \(R\in\mathcal P\) be a declared
non-auxiliary relationship type with \(\alpha(R)\geq 2\), and let \(I\)
be a finite instance over the non-auxiliary KG-ER\(_0\) vocabulary of
\(\mathcal S\).  Put
\(A=\operatorname{adom}(I)\), and let
\(\mathcal S_{\mathsf{nk}}\) be the non-key part of \(\mathcal S\), i.e.,
all constraints except key dependencies; the excluded keys include the
explicit keys on the auxiliary relationships \(P_\iota\).
In every non-failing chase branch of \(I\) with
\(\mathcal S_{\mathsf{nk}}\), with fresh existential witnesses and cover
constraints treated by case splitting, the restriction of \(R\) to the
initial active domain is unchanged: if an \(R\)-fact
\(R(\mathbf a)\) occurs in the branch with
\(\mathbf a\in A^{\alpha(R)}\), then \(R(\mathbf a)\in I\).
Equivalently, every \(R\)-fact in the branch is either an original
\(R\)-fact of \(I\), or has at least one component outside \(A\).

Thus, when the chase is viewed in two stages, the non-key stage can add
only \(R\)-facts containing fresh witnesses.  The subsequent key stage
does not prune tuples; it only identifies terms, or fails.  Consequently,
any new \(R\)-fact entirely over \(A\) in the full chase must arise from
a key equality that identifies fresh witnesses with elements of \(A\).
\end{lemma}

\begin{proof}

\smallskip\noindent\textbf{Strategy.}
The proof is by induction on the non-key chase steps.  The induction
hypothesis is a \emph{simultaneous invariant} maintaining two facts at
once.  One fact suffices to state the conclusion for \(R\), but the
second fact, about auxiliary \(P_\iota\)-facts, is needed to make the
induction go through, because the Id-generated reverse inclusion can
create a new \(R\)-fact from an auxiliary fact.

\smallskip\noindent\textbf{The simultaneous invariant.}
The invariant has two clauses, both required to hold throughout the
chase.

\emph{Clause (i) --- about \(R\)-facts.}
Each \(R\)-fact in the current instance is either (a) already present
in \(I\), or (b) contains at least one fresh witness.

\emph{Clause (ii) --- origin form of auxiliary facts.}
Every auxiliary fact \(P_\iota(c,a_1,\ldots,a_m)\) has one of two
specific origins:
\begin{itemize}
\item \emph{Origin type 1}: it was generated by the inclusion
\(E\subseteq P_\iota.1\).  In this case all of \(a_1,\ldots,a_m\) are
fresh witnesses introduced at the moment of creation.
\item \emph{Origin type 2}: it was generated by the forward
Id-inclusion from some identifying fact \(P_j(c,a_j)\).  In this case
the first component \(c\) and the \((j+1)\)-st component \(a_j\) are
copied from the source fact \(P_j(c,a_j)\); all other non-first
components (i.e., positions \(k+1\) for \(k\neq j\)) are fresh witnesses
introduced at the moment of creation.
\end{itemize}

This precise structural description of auxiliary facts is the key
insight: every \(P_\iota\)-fact has \emph{at most} one ``old''
non-first component (and even that one is old only if the source
\(P_j\)-fact had an old component there).  All other non-first
positions are fresh.

\smallskip\noindent\textbf{Base case.}
Initially the instance is \(I\), which contains no auxiliary facts, so
clause (ii) is vacuously satisfied; and every \(R\)-fact in \(I\) is,
by definition, present in \(I\), so clause (i) holds.

\smallskip\noindent\textbf{Trivial steps preserve the invariant.}
The following chase steps have nothing to do with \(R\)-facts or
auxiliaries:
\begin{itemize}
\item \emph{Type declarations} (typing inclusions from
\(\mathcal B_{\mathcal S}\)): only add unary facts \(T(c)\).
\item \emph{Entity inclusions} (IsA): only add unary facts.
\item \emph{Cover constraints}: cause a branch split, not new facts.
\item \emph{Disjointness constraints}: either do not fire or close the
branch.
\end{itemize}
None of these touches \(R\)-facts or creates \(P_\iota\)-facts, so both
clauses are preserved unchanged.

\smallskip\noindent\textbf{Mandatory inclusions \(E\subseteq R.i\) preserve the invariant.}
A mandatory inclusion \(E\subseteq R.i\) for a non-auxiliary
relationship \(R\) (this comes from \(\mathcal N_{\mathcal S}\)) is an
existential rule: whenever \(E(c)\) holds and no \(R\)-tuple has \(c\)
at position \(i\), introduce one.  The new \(R\)-tuple has \(c\) at
position \(i\) and fresh existentially witnessed values at the other
\(\alpha(R)-1\) positions.  Since \(\alpha(R)\geq 2\), the new tuple
has at least one fresh component, so clause (i) holds for the new
\(R\)-fact (it has a fresh witness, even though \(c\) may be old).

Mandatory inclusions for other relationships \(Q\neq R\) do not affect
\(R\)-facts at all, so clause (i) for \(R\) is trivially preserved.
Mandatory inclusions do not generate \(P_\iota\)-facts, so clause (ii)
is also untouched.

\smallskip\noindent\textbf{Id-generated forward inclusion preserves the invariant.}
This is the rule
\[
P_j[\ols{i_j},i_j]\subseteq P_\iota[1,j+1].
\]
It fires when there is a \(P_j\)-tuple in positions
\((\ols{i_j},i_j)\) --- say with values \((c,a_j)\) --- and no
corresponding \(P_\iota\)-tuple yet.  It creates a new
\(P_\iota\)-fact with \(c\) at position 1, \(a_j\) at position
\(j+1\), and fresh existentially witnessed values at all other
non-first positions.

This generated \(P_\iota\)-fact satisfies clause (ii)'s origin type 2
with the source \(P_j(c,a_j)\).  The clause-(ii) form is recorded
exactly as required: position 1 and position \(j+1\) are copies from
the source; positions \(k+1\) for \(k\neq j\) are fresh witnesses.

Crucially, this holds \emph{regardless of whether the source
\(P_j(c,a_j)\) is an original fact of \(I\) or was generated earlier
by some chase step}.  The invariant is purely about the \emph{form}
of the auxiliary fact, not about the provenance of its source.

Clause (i) is unaffected because the forward inclusion only generates
auxiliary facts, not \(R\)-facts.

\smallskip\noindent\textbf{Id-generated reverse inclusion preserves the invariant.}
This is the dangerous step:
\[
P_\iota[1,j+1]\subseteq P_j[\ols{i_j},i_j].
\]
It fires when there is a \(P_\iota\)-tuple with values at positions 1
and \(j+1\) --- say \((c,a)\) --- and no corresponding \(P_j\)-tuple
yet.  It creates a new \(P_j\)-fact with \(c\) at position
\(\ols{i_j}\) and \(a\) at position \(i_j\).

If \(P_j=R\), this could potentially create a new \(R\)-fact,
threatening clause (i).  We case-split on the origin of the source
\(P_\iota\)-fact.

\emph{Case A --- source has origin type 1} (generated by
\(E\subseteq P_\iota.1\)).
By the invariant, the \((j+1)\)-st component of the source is fresh.
So the value \(a\) being projected to position \(i_j\) of the new
\(P_j\)-fact is fresh.  The new \(P_j\)-fact thus contains a fresh
witness.  Clause (i) preserved.

\emph{Case B --- source has origin type 2} (generated from some
identifying fact \(P_k(c,a_k)\)).
Two sub-cases on the relationship between \(j\) (the index governing
the reverse inclusion currently firing) and \(k\) (the index that
generated the source \(P_\iota\)-fact).

\emph{Sub-case B1: \(j=k\).}
The reverse inclusion projects positions \((1,j+1)\) of the source
\(P_\iota\)-fact back to \(P_j\).  By origin type 2 with \(k=j\),
these positions hold exactly the source \(P_j(c,a_j)\) data.  So the
reverse inclusion \emph{reproduces the same \(P_j\)-fact}
\(P_j(c,a_j)\) that was the original source.

Now: if this is an \(R\)-fact (i.e., \(P_j=R\)) entirely over \(A\),
then by the induction hypothesis (clause (i) applied to the source)
the source fact was already in \(I\).  So the reproduced fact is
already in \(I\).  No new \(R\)-fact over \(A\) is created.
If the source had a fresh component, the reproduced fact also has a
fresh component.  Clause (i) preserved.

\emph{Sub-case B2: \(j\neq k\).}
The \((j+1)\)-st component of the source \(P_\iota\)-fact is, by
origin type 2 with \(k\neq j\), one of the fresh witnesses introduced
when the source \(P_\iota\)-fact was created.  So the value \(a\)
being projected to position \(i_j\) of the new \(P_j\)-fact is fresh.
The new \(P_j\)-fact contains a fresh witness.  Clause (i) preserved.

In all sub-cases of Case B, either the reverse inclusion reproduces a
\(P_j\)-fact already covered by clause (i), or it creates a
\(P_j\)-fact with a fresh witness.

Clause (ii) is unaffected because the reverse inclusion creates a
non-auxiliary \(P_j\)-fact, not a \(P_\iota\)-fact.

\smallskip\noindent\textbf{Conclusion of the inductive step.}
Every non-key chase rule has been checked: background and type rules,
mandatory existentials, forward Id-inclusions, and reverse
Id-inclusions.  All preserve the simultaneous invariant.  By
induction, the invariant holds at every step of every non-failing
non-key branch.

In particular, clause (i) gives the first conclusion of the theorem:
no non-key chase step ever creates a new \(R\)-fact entirely over
\(A\).  Every all-old \(R\)-fact in any non-key branch was already
in \(I\).

\smallskip\noindent\textbf{Final consequence about keys.}
Enforcing keys is equality-generating, not fact-generating.  A key
step either identifies two terms (merging equivalence classes) or
fails (if it requires impossible equality).  It does not add or
delete facts.

So if the full chase (with keys) produces a new \(R\)-fact entirely
over \(A\), the only way this can happen is: some \(R\)-fact with a
fresh witness was already in the non-key chase, and a key step then
identified that fresh witness with an element of \(A\), turning the
fact into an all-\(A\) tuple.

This is the separation: fresh witnesses can be introduced only by
non-key steps, and they can be merged into \(A\) only by key steps.
The two roles are cleanly disjoint.
\end{proof}

\begin{lemma}[\colorbox{yellow}{Old-domain equality preservation}]
\label{lem:kger-old-domain-equality}
Let \(\mathcal S\) be the compiled KG-ER\(_0^+\) expansion of an
ER-satisfiable KG-ER\(_0\) schema, and let \(I\) be a finite instance
over the non-auxiliary KG-ER\(_0\) vocabulary of \(\mathcal S\).  Put
\(A=\operatorname{adom}(I)\).  Fix a non-failing branch of the chase of
\(I\) with
\(\mathcal B_{\mathcal S}\cup\mathcal D_{\mathcal S}\cup
\mathcal N_{\mathcal S}\), before applying keys.  Let \(I_D\) be the
instance obtained from the same initial instance by chasing only with
\(\mathcal B_{\mathcal S}\cup\mathcal D_{\mathcal S}\), and let
\(I_{DN}\) be the instance obtained on the fixed branch by also chasing
\(\mathcal N_{\mathcal S}\), again before applying keys.

Let \(\equiv_D\) be the equivalence relation on terms generated by the
key chase of \(I_D\) with \(\mathcal K_{\mathcal S}\), and let
\(\equiv_{DN}\) be the equivalence relation generated by the key chase of
\(I_{DN}\) with \(\mathcal K_{\mathcal S}\).  Assume that the key chase
of \(I_{DN}\) does not fail; since \(I_D\subseteq I_{DN}\) and key
dependencies are preserved under taking sub-instances, the key chase of
\(I_D\) does not fail either.  Then
\[
a\equiv_{DN} b \quad\text{iff}\quad a\equiv_D b
\qquad\text{for all }a,b\in A .
\]
Equivalently, adding the remaining non-key constraints before the key
chase does not create any new equality between initial active-domain
terms.
\end{lemma}

\begin{proof}

\smallskip\noindent\textbf{Step 1 --- Easy direction.}
If \(a\equiv_D b\), then \(a\equiv_{DN} b\) by monotonicity:
\(I_D\subseteq I_{DN}\), so every key step available in the
\(I_D\)-chase is also available in the \(I_{DN}\)-chase, and
equivalences accumulate.

\smallskip\noindent\textbf{Step 2 --- Hard direction: a first-bad-step
argument.}
Suppose for contradiction that some \(a,b\in A\) satisfy
\(a\equiv_{DN} b\) but \(a\not\equiv_D b\).  Then the \(I_{DN}\) key
chase generates a new \(A\)-\(A\) equality not present in \(\equiv_D\).
Pick the first key step (in the chase order) that creates such a new
equality; call it the \emph{chosen first step}.  We derive a
contradiction by showing this step cannot in fact be the first
creator of a new \(A\)-\(A\) equality.

\smallskip\noindent\textbf{Step 3 --- Terminology: \(D\)-supported facts.}
A fact is \emph{\(D\)-supported} if it belongs to \(I_D\).  We need to
know what the non-\(D\)-supported facts of \(I_{DN}\) look like ---
i.e., what is in \(I_{DN}\) but not in \(I_D\).

By Lemma~\ref{lem:kger-key-nonkey-separation}, applied not just to a
single \(R\) but to all non-auxiliary relationships of arity at least
two simultaneously (Lemma~\ref{lem:kger-key-nonkey-separation}'s proof
maintains its invariant globally across all chase steps and all
relationships, even though its statement names a single \(R\)), every
non-\(D\)-supported relationship fact of arity at least two in
\(I_{DN}\) contains at least one term not in \(I_D\).  Such a term is
called a \emph{fresh witness}.  These fresh witnesses come from one of
two sources:
\begin{itemize}
\item a rule in \(\mathcal N_{\mathcal S}\) (e.g., a mandatory
existential like \(E\subseteq R.i\) fires and introduces a new
witness), or
\item a rule in \(\mathcal D_{\mathcal S}\) fired from a
non-\(D\)-supported trigger (e.g., a forward Id-inclusion firing on a
\(P_j\)-fact that itself is not in \(I_D\)).
\end{itemize}
Unary type facts \(T(c)\) are also added by \(\mathcal N_{\mathcal S}\),
but unary facts carry no key dependencies and are irrelevant to the key
chase.

For auxiliary \(P_\iota\)-facts the same kind of invariant applies: the
Id-inclusions either copy values from binary projections or introduce
existentials, and any positions filled by fresh existentials sit
outside \(I_D\) when the trigger is itself outside \(I_D\).

\smallskip\noindent\textbf{Step 4 --- Trivial sub-case: both facts
\(D\)-supported.}
The chosen first step uses two facts of the same relationship, agreeing
on the determinant of a declared key.  If both facts are
\(D\)-supported (both in \(I_D\)), then the same key step is also
available in the \(I_D\) key chase.  So the equality it generates is
already in \(\equiv_D\), contradicting the assumption that it creates a
new equality.  This sub-case is done.

So from here on, at least one of the two facts in the chosen step is
not \(D\)-supported.

\smallskip\noindent\textbf{Step 5 --- Setting up old traces via
minimality.}
Define the \emph{old trace} of a term \(t\) at any point in the key
chase as the set of \(A\)-elements equivalent to \(t\) at that point.
By minimality of the chosen first step:
\begin{quote}
Before the chosen step, no current equivalence class contains
representatives of two distinct \(\equiv_D\)-classes from \(A\).
\end{quote}
(If it did, that earlier merger would have been an earlier ``first
bad step'' --- contradicting minimality.)  Therefore every term in the
chase, just before the chosen step, has an \emph{unambiguous} old
trace: either no \(A\)-representative at all, or exactly one
\(\equiv_D\)-class worth of \(A\)-representatives.

This is the inductive use of minimality.  From now on we can speak of
``the old trace of \(t\)'' unambiguously, identifying it with a single
\(\equiv_D\)-class (or ``none'').

\smallskip\noindent\textbf{Step 6 --- Component classification.}
Use the origin invariant of
Lemma~\ref{lem:kger-key-nonkey-separation}.  Each component of a
non-\(D\)-supported fact falls into one of two categories:
\begin{itemize}
\item \emph{Type (a): copied from a \(D\)-supported fact.}  Two
situations: (a1) the component occurs in an initial unary type fact of
\(I\) --- so it is an \(A\)-element with its own identity; (a2) the
component was copied through a \(D\)-supported binary projection (e.g.,
position 1 or position \(j+1\) of a \(P_\iota\)-fact whose source is
in \(I_D\)).
\item \emph{Type (b): a fresh witness.}  Either introduced by an
\(\mathcal N_{\mathcal S}\)-rule, or by a \(\mathcal D_{\mathcal S}\)-rule
fired from a non-\(D\)-supported trigger.
\end{itemize}
We now analyse how each kind of component contributes to determinant
agreement in the chosen step.

\smallskip\noindent\textbf{Step 7 --- Type-(b) components cannot start
the first bad merger.}
A type-(b) component does not occur in \(I_D\), so it has no old trace
at creation time.  The only way for it to acquire an old trace before
the chosen step is via an earlier key step.  By minimality (Step~5),
any such earlier step can have copied into it at most a single
\(\equiv_D\)-class of \(A\).  So a type-(b) component's current old
trace is either empty or a single \(\equiv_D\)-class.

Hence determinant agreement involving a type-(b) component cannot be
the first occasion on which two distinct old traces are merged: it
carries at most one old trace, so even when the step fires and
propagates this trace to the conclusion side, the result is just
copying that single trace --- not merging two.

\smallskip\noindent\textbf{Step 8 --- Type-(a) components, determinant
agreement.}
This is where the actual content lies.  The first key step fires
because two facts \(\phi_1,\phi_2\) of the same relationship agree on
the determinant \(X\) of a declared key.  Each determinant position
\(x\in X\) has a component in \(\phi_1\) (call it \(t\)) and a
component in \(\phi_2\) (call it \(t'\)), with \(t\equiv_{DN} t'\)
before the step.

If either \(t\) or \(t'\) is type-(b), Step~7 already shows this
position cannot drive the first \(A\)-\(A\) merger.  So suppose both
\(t\) and \(t'\) are type-(a).

\emph{Sub-case (a1): copied from an initial unary fact.}  Then \(t\)
is literally an \(A\)-element with its identity fixed from the start.
The old trace of \(t\) is just \(\{t\}\) as an \(\equiv_D\)-class (or
its \(\equiv_D\)-class if any earlier \(\equiv_D\)-equality has merged
it).  Same for \(t'\).  Agreement on this determinant position can only
be agreement between two old traces that are already equal --- either
literally equal \(A\)-elements or already in the same
\(\equiv_D\)-class.

\emph{Sub-case (a2): copied from a \(D\)-supported binary projection.}
Now \(\phi_1\) has its determinant component \(t\) copied from some
\(D\)-supported fact \(\phi_1^D\in I_D\), and similarly \(\phi_2\)'s
\(t'\) is copied from \(\phi_2^D\in I_D\).  Since the copying preserves
the position's value through \(\mathcal D_{\mathcal S}\)-rules, the same
determinant agreement \(t\equiv t'\) is realised at the
\(D\)-supported level: \(\phi_1^D\) and \(\phi_2^D\) have the same
components at corresponding positions, modulo \(\equiv_D\).  Therefore
the key chase of \(I_D\) already fires the same key step on
\(\phi_1^D\) and \(\phi_2^D\), generating the same equality between the
corresponding old traces.

In both type-(a) sub-cases, the current step is merging two old traces
only when those traces are already the same \(\equiv_D\)-class.

\smallskip\noindent\textbf{Step 9 --- Auxiliary keys.}
The argument above is enough for both auxiliary keys on \(P_\iota\),
because both have determinants made of type-(a)/(b) components in the
structure provided by Lemma~\ref{lem:kger-key-nonkey-separation}.

\emph{Auxiliary key \(P_\iota.1\to P_\iota.2\cdots P_\iota.(m+1)\):}
determinant is position 1.  By
Lemma~\ref{lem:kger-key-nonkey-separation}'s invariant on
\(P_\iota\)-facts, position 1 of any \(P_\iota\)-fact is copied from a
\(D\)-supported identifying projection (the forward inclusion always
copies into position 1 from a source \(P_k\)-fact).  So the determinant
component is type-(a2), and any old trace propagated to a non-first
position by this key step is one already forced by the corresponding
\(D\)-supported projection --- already in \(\equiv_D\).

\emph{Auxiliary key \(P_\iota.2\cdots P_\iota.(m+1)\to P_\iota.1\):}
determinant is the bundle of all identifying positions.  Two
\(P_\iota\)-facts fire this key only when they agree on every
identifying component.  For positions whose value has an old trace,
the agreement is again type-(a2), so already in \(\equiv_D\).  For
positions where one fact's value has no old trace (type-(b) component),
Step~7 applies: that position carries at most one old trace, so the
step cannot be the first to merge two old traces.

\smallskip\noindent\textbf{Step 10 --- Contradiction.}
The chosen first step either:
\begin{itemize}
\item merges a class with no old trace into a class with at most one
old trace (Step~7 covers this --- no new \(A\)-\(A\) equality), or
\item merges two classes whose old traces are already equal modulo
\(\equiv_D\) (Step~8 covers this --- also not a new \(A\)-\(A\)
equality).
\end{itemize}
In neither case does the step actually create a new equality between
two distinct \(\equiv_D\)-classes from \(A\).  This contradicts the
choice of the step as ``the first one that does so.''  Therefore no
such step exists, and the hard direction is proved.
\end{proof}

\begin{lemma}[\colorbox{yellow}{Old-domain binary-fact preservation}]
\label{lem:kger-old-domain-binary-fact}
Let \(\mathcal S\), \(I\), \(A\), \(I_D\), \(I_{DN}\),
\(\equiv_D\), and \(\equiv_{DN}\) be as in
Lemma~\ref{lem:kger-old-domain-equality}.  Let \(R\in\mathcal P\) be a
non-auxiliary binary relationship.  Suppose that the key chase of
\(I_{DN}\) does not fail.  If, in the quotient of \(I_{DN}\) by
\(\equiv_{DN}\), there is an \(R\)-fact with ordered old-domain trace
\((a,b)\in A^2\), then, in the quotient of \(I_D\) by \(\equiv_D\),
there is already an \(R\)-fact with ordered old-domain trace \((a,b)\).
\end{lemma}

\begin{proof}

\smallskip\noindent\textbf{Step 1 --- First-bad-trace setup.}
Suppose for contradiction that the conclusion fails.  Then there is
some ordered pair \((a,b)\in A^2\) such that the \(I_{DN}\)-quotient
contains an \(R\)-fact with trace \((a,b)\) but the \(I_D\)-quotient
does not.

Consider the first key step in the \(I_{DN}\) chase that creates a
witnessing \(R\)-fact of this kind --- i.e., the first key step after
which the \(I_{DN}\)-quotient contains an \(R\)-fact with an ordered
\(A^2\)-trace that the \(I_D\)-quotient has never realised.  Call this
the \emph{chosen first bad step}, and call the offending \(R\)-fact
\(R(s,t)\) with \(\equiv_{DN}\)-trace \((a,b)\).

We will derive a contradiction by showing the chosen step cannot in
fact be the first creator of such a trace.

\smallskip\noindent\textbf{Step 2 --- Case 1: the witnessing
\(R\)-fact is \(D\)-supported.}
Suppose \(R(s,t)\in I_D\).  Then the only way it acquires the trace
\((a,b)\) in the \(I_{DN}\)-quotient is via equalities
\(s\equiv_{DN} a\) and \(t\equiv_{DN} b\).

By Lemma~\ref{lem:kger-old-domain-equality}, the restriction of
\(\equiv_{DN}\) to \(A\) equals the restriction of \(\equiv_D\) to
\(A\).  So if \(s,t\in A\), then \(s\equiv_{DN} a\) gives
\(s\equiv_D a\), and similarly for \(t\).  But then \(R(s,t)\in I_D\)
has trace \((a,b)\) already in \(\equiv_D\), contradicting the choice
of \((a,b)\) as not realised in the \(I_D\)-quotient.

If \(s\) or \(t\) is a \(\mathcal D_{\mathcal S}\)-introduced fresh
witness in \(\operatorname{adom}(I_D)\setminus A\), the conclusion is
the same: any equalities propagating \(s\) to \(a\) and \(t\) to \(b\)
in \(\equiv_{DN}\) are already present in \(\equiv_D\), by the same
first-bad-trace argument as in
Lemma~\ref{lem:kger-old-domain-equality} applied to these terms.
Hence \(R(s,t)\) in the \(I_D\)-quotient has trace \((a,b)\) ---
contradiction.

So the witnessing fact \(R(s,t)\) cannot be \(D\)-supported.

\smallskip\noindent\textbf{Step 3 --- Case 2: the witnessing
\(R\)-fact is non-\(D\)-supported.}
Then \(R(s,t)\in I_{DN}\setminus I_D\).  By
Lemma~\ref{lem:kger-key-nonkey-separation} applied at the \(I_{DN}\)
level, every \(R\)-fact in \(I_{DN}\) is either in \(I\) (hence in
\(I_D\)) or has at least one fresh witness --- a term not in
\(\operatorname{adom}(I_D)\).  So at least one of \(s,t\) is a fresh
witness, introduced by an \(\mathcal N_{\mathcal S}\)-rule or by a
\(\mathcal D_{\mathcal S}\)-rule fired from a non-\(D\)-supported
trigger.

\smallskip\noindent\textbf{Step 4 --- Component classification
(mirroring Lemma~\ref{lem:kger-old-domain-equality}).}
Reuse the component classification from
Lemma~\ref{lem:kger-old-domain-equality}.  Each component of the
non-\(D\)-supported fact \(R(s,t)\) is either:
\begin{itemize}
\item \emph{Type (a):} copied from \(D\)-supported material (either an
initial unary type fact giving an \(A\)-element directly, or a
\(D\)-supported binary projection).
\item \emph{Type (b):} a fresh witness, introduced by an
\(\mathcal N_{\mathcal S}\)-rule or by a \(\mathcal D_{\mathcal S}\)-rule
fired from a non-\(D\)-supported trigger.
\end{itemize}
Since \(R(s,t)\) is non-\(D\)-supported, at least one of \(s,t\) is
type-(b).

\smallskip\noindent\textbf{Step 5 --- Old traces on type-(a) and
type-(b) components.}
For each component of \(R(s,t)\) that has an old trace in
\(\equiv_{DN}\):
\begin{itemize}
\item \emph{Type-(a) component:} Its old trace is already visible at
the \(D\)-level.  If it came from an initial unary fact, it is
literally an \(A\)-element.  If it came from a \(D\)-supported binary
projection, the same projection exists in \(I_D\), so the corresponding
component's old trace is realised in \(\equiv_D\) via the same chain
of \(D\)-supported facts.
\item \emph{Type-(b) component:} Has no old trace at creation.  Any
old trace it later acquires must come from earlier key steps.  By the
minimality of the chosen first bad step (Step~1), each earlier key
step has effects already visible at the \(D\)-level --- otherwise that
earlier step would have been a first bad step.  So a type-(b)
component's current old trace, if any, corresponds to an analogous
identification in \(\equiv_D\) between corresponding \(D\)-supported
terms.
\end{itemize}

\smallskip\noindent\textbf{Step 6 --- Putting the trace together.}
By Step~5, the assignments \(s\mapsto a\) and \(t\mapsto b\) (the old
traces of \(R(s,t)\)'s components in \(\equiv_{DN}\)) each correspond
to identifications already realised at the \(D\)-level.  The ordered
pair \((a,b)\) of old traces is therefore reflected, position by
position, in an \(R\)-fact at the \(D\)-level: there exists some
\(D\)-supported \(R\)-fact \(R(s',t')\) (possibly with \(s',t'\) being
\(\mathcal D_{\mathcal S}\)-introduced fresh witnesses, possibly
\(A\)-elements) such that \(s'\equiv_D a\) and \(t'\equiv_D b\).

The key point is that ordered binary projections through
\(\mathcal D_{\mathcal S}\)-rules are determined by the
\(D\)-supported source facts: the inclusion
\(P_j[\ols{i_j},i_j]\subseteq P_\iota[1,j+1]\) sends position
\(\ols{i_j}\) to position 1, etc.  So a \(D\)-supported source
\(R\)-fact projects to a specific ordering in the auxiliary \(P_\iota\),
and vice versa.  Component ordering is preserved through chained
projections.

\smallskip\noindent\textbf{Step 7 --- Contradiction.}
Step~6 shows that the ordered trace \((a,b)\) is realised by some
\(R\)-fact in the \(I_D\)-quotient --- contradicting the choice of
\((a,b)\) at Step~1 as a trace not realised in the \(I_D\)-quotient.
Therefore no chosen first bad step exists, and the lemma holds.
\end{proof}

\noindent\textbf{Statement of Theorem~\ref{thm:kger-pure-key-separation} (\colorbox{yellow}{Relative key/non-key entailment separation}).}
Let \(\mathcal S\) be the compiled KG-ER\(_0^+\) expansion of an
ER-satisfiable KG-ER\(_0\) schema.  Then, for every key dependency
\(\kappa\) over a non-auxiliary
relationship \(R\in\mathcal P\),
\[
\mathcal B_{\mathcal S}\cup\mathcal D_{\mathcal S}\cup
\mathcal K_{\mathcal S}\cup\mathcal N_{\mathcal S}\models\kappa
\quad\text{iff}\quad
\mathcal B_{\mathcal S}\cup\mathcal D_{\mathcal S}\cup
\mathcal K_{\mathcal S}\models\kappa .
\]
Thus, once the inclusion dependencies introduced by the KG-ER\(_0^+\)
encoding of identification constraints are fixed, the remaining non-key
constraints do not contribute to the entailment of keys.

\begin{proof}[Proof of Theorem~\ref{thm:kger-pure-key-separation} (Relative key/non-key entailment separation)]
The right-to-left implication is immediate by monotonicity of entailment.
For the converse, let \(\kappa\) be a key dependency over
\(R\in\mathcal P\).  By the vocabulary convention,
\(\alpha(R)\geq2\).

Use the standard chase test for implication of the key \(\kappa\).  Start
from the two-row \(R\)-tableau \(T_\kappa\) whose rows agree exactly on
the determinant of \(\kappa\), together with the type facts required by
\(\mathcal B_{\mathcal S}\).  Let \(A\) be the active domain of this
initial tableau.
Since \(\mathcal S\not\models R=\emptyset\), there is a model
\(M\models\mathcal S\) with \(R^M\neq\emptyset\).  Mapping both tableau
rows to the same tuple of \(R^M\) gives a homomorphism into a model of
\(\mathcal S\); hence the non-key chase has a non-failing branch, with
cover choices made according to \(M\), and the subsequent key chase on
that branch does not fail.

Let \(I_D\) be the result of chasing \(T_\kappa\) with
\(\mathcal B_{\mathcal S}\cup\mathcal D_{\mathcal S}\) along this branch,
before applying keys, and let \(I_{DN}\) be the corresponding result
when \(\mathcal N_{\mathcal S}\) is also chased before applying keys.
By Lemma~\ref{lem:kger-old-domain-equality}, the restriction to \(A\) of
any key quotient obtained from \(I_{DN}\) is already generated by
\(\mathcal B_{\mathcal S}\cup\mathcal D_{\mathcal S}\cup
\mathcal K_{\mathcal S}\).  Therefore, if
\(\mathcal S\models\kappa\), the chase of \(T_\kappa\) with
\(\mathcal B_{\mathcal S}\cup\mathcal D_{\mathcal S}\cup
\mathcal K_{\mathcal S}\) already identifies the corresponding
components required by \(\kappa\).  By soundness and completeness of the
chase for key implication,
\[
\mathcal B_{\mathcal S}\cup\mathcal D_{\mathcal S}\cup
\mathcal K_{\mathcal S}\models\kappa .
\]
This proves the converse implication.
\end{proof}

\noindent\textbf{Statement of Theorem~\ref{thm:kger-5} (\colorbox{yellow}{5$^{th}$ normal form of KG-ER$_0$}).}
Let \(\mathcal S\) be an ER-satisfiable KG-ER\(_0\) schema, and let
\(R\) be a declared relationship type of arity \(n\).  Let
\(\Sigma_R\) be the set of all key dependencies
\(X\to R.1\cdots R.n\) entailed by \(\mathcal S\), with
\(X\subseteq\{R.1,\ldots,R.n\}\).  Then, for every join
dependency \(J\) over the attributes of \(R\),
\[
\mathcal S\models J
\quad\text{iff}\quad
\Sigma_R\models J .
\]
Hence every such relationship type \(R\) is in 5$^{th}$ normal form.

\begin{proof}[Proof of Theorem~\ref{thm:kger-5} (5$^{th}$ normal form of KG-ER$_0$)]

\smallskip\noindent\textbf{Step 1 --- Reduce to the compiled schema.}
By Theorem~\ref{thm:kger-definitional-extension} (KG-ER\(_0^+\) as a
conservative extension), entailments over the original vocabulary are
preserved and reflected by passing to the compiled expansion
\(\mathcal S^+\).  Since both \(\mathcal S\models J\) and
\(\Sigma_R\models J\) are statements over \(R\)'s vocabulary, we may
work in \(\mathcal S^+\) instead of \(\mathcal S\) --- gaining the
explicit auxiliaries \(P_\iota\) and the ordinary key/inclusion
machinery in \(\mathcal D_{\mathcal S}\), \(\mathcal K_{\mathcal S}\).
The original Id constraints are now replaced by standard keys and
inclusions, so the chase becomes a standard chase.

\smallskip\noindent\textbf{Step 2 --- Easy direction (\(\Leftarrow\)).}
If \(\Sigma_R\models J\), then \(\mathcal S\models J\) by transitivity:
every dependency in \(\Sigma_R\) is entailed by \(\mathcal S\) by
definition, so every JD implied by \(\Sigma_R\) is also entailed by
\(\mathcal S\).

\smallskip\noindent\textbf{Step 3 --- Hard direction
(\(\Rightarrow\)): chase test setup.}
Use the standard chase test for JD implication.  For
\(J=\bowtie[U_1,\ldots,U_k]\), the \emph{JD tableau} is a relation
containing \(k\) rows, each a copy of \(R\)'s schema.  Distinguished
variables (call them \(\overline{x}\)) appear at positions in \(U_i\)
in row \(i\); fresh existentially-bound variables fill all other
positions.  The chase forces \(J\) if and only if some row eventually
equals the distinguished tuple \(\overline{x}\) after the chase.

Augment the tableau with the type facts required by \(R\)'s
declaration (to ensure typing is satisfied).  Let \(A\) be the active
domain of this initial tableau --- the distinguished variables plus
the fresh variables across all rows.

Since \(R\in\mathcal P\), the vocabulary convention gives \(n\geq 2\)
(so the JD test is non-degenerate).

\smallskip\noindent\textbf{Step 4 --- A non-failing chase branch
exists.}
Since \(\mathcal S\) is ER-satisfiable, there is a model
\(M\models\mathcal S\) with \(R^M\neq\emptyset\).  Pick any tuple
\(\overline{t}\in R^M\) and send every row of the JD tableau to
\(\overline{t}\): this is a homomorphism into \(M\) because each row,
projected onto its \(U_i\), agrees with \(\overline{x}\), and
\(\overline{x}\) maps to \(\overline{t}\).  The type facts of the
tableau map to type facts true in \(M\) by typing.

Having a homomorphism into a real model \(M\) lets us pick branches
consistent with \(M\) at every step.  Along such a branch: existential
rules pick witnesses according to \(M\); disjointness constraints do
not fire (they hold in \(M\)); the subsequent key chase does not fail
(keys hold in \(M\)).  So at least one non-failing branch exists; the
proof works along one such branch.

\smallskip\noindent\textbf{Step 5 --- Restrict to \(A\)-equalities via
Lemma~\ref{lem:kger-key-nonkey-separation}.}
By Lemma~\ref{lem:kger-key-nonkey-separation}, along the non-failing
branch the non-key part of the chase cannot create a new \(R\)-fact
whose components all lie in \(A\).  Every \(R\)-fact produced by
non-key rules has at least one fresh witness outside \(A\).

\emph{Consequence.}  If the chase forces the distinguished tuple
\(\overline{x}\), the only mechanism is quotienting tableau variables
by key-step equalities, collapsing rows onto each other until one row
becomes the distinguished tuple.  So the success of the chase is
entirely determined by the equivalence relation on \(A\) generated by
key steps along the branch.

\emph{Reduction of the proof obligation.}  We need to show: every
\(A\)-\(A\) equality generated by a key step is a consequence of
\(\Sigma_R\).  If so, the chase using \(\Sigma_R\) alone produces the
same \(A\)-equivalences, hence the same identifications of rows, hence
forces the distinguished tuple.  By soundness and completeness of the
chase for JD implication, this gives \(\Sigma_R\models J\).

So everything reduces to a case analysis on the relationship symbol of
each key step.

\smallskip\noindent\textbf{Step 6 --- Case A: key step on \(R\)
itself.}
The step uses a key dependency \(X\to R.1\cdots R.n\) on \(R\).  By
definition this dependency is entailed by \(\mathcal S\), hence in
\(\Sigma_R\).  The \(A\)-\(A\) equality it generates is therefore a
direct consequence of a member of \(\Sigma_R\).

\smallskip\noindent\textbf{Step 7 --- Case B: key step on an auxiliary
\(P_\iota\).}
This is the substantive case.  Let
\(\iota=\textsc{Id}(E\ (P_1.i_1,\ldots,P_m.i_m))\) be the
identification constraint associated with \(P_\iota\).

\emph{Reduction to binary \(R\).}  By the definition of \(\textsc{Id}\),
every \(P_k\) in \(\iota\) has \(\alpha(P_k)=2\).  A key step on
\(P_\iota\) can only propagate to a non-auxiliary \(R\) through the
bidirectional inclusions \(P_\iota[1,k+1]=P_k[\ols{i_k},i_k]\).  If
\(R\) has arity greater than two, no Id constraint involves \(R\) as
an identifying relationship, so no \(P_\iota\)-key can affect
\(R\)-variables --- done.  So assume \(R=P_j\) for some \(j\), with
\(R\) binary.

The two auxiliary keys on \(P_\iota\) are:
\begin{itemize}
\item \emph{First key:} \(P_\iota.1\to P_\iota.2\cdots P_\iota.(m+1)\)
--- ``the entity determines its identifying tuple''.
\item \emph{Second key:}
\(P_\iota.2\cdots P_\iota.(m+1)\to P_\iota.1\) --- ``the identifying
tuple determines the entity''.
\end{itemize}
We analyse each.

\smallskip\noindent\textbf{Step 7a --- First auxiliary key.}
Restrict the first key \(P_\iota.1\to P_\iota.2\cdots P_\iota.(m+1)\)
to positions \(\{1,j+1\}\): it says any two \(P_\iota\)-tuples
agreeing on position 1 must agree on position \(j+1\).  Via the
bidirectional inclusion \(P_\iota[1,j+1]=P_j[\ols{i_j},i_j]\), this
translates to: any two \(P_j\)-tuples agreeing on position
\(\ols{i_j}\) must agree on position \(i_j\).  That is the functional
dependency
\[
\ols{i_j}\to i_j \quad\text{on } P_j.
\]

Since \(P_j\) is binary with positions
\(\{\ols{i_j},i_j\}=\{1,2\}\), this FD together with the trivial
\(\ols{i_j}\to\ols{i_j}\) is the full key
\[
\ols{i_j}\to P_j.1\,P_j.2
\]
on \(P_j\).  This is entailed by \(\mathcal S\) (since the key on
\(P_\iota\) is in \(\mathcal K_{\mathcal S}\) and the bidirectional
inclusion is in \(\mathcal D_{\mathcal S}\)), so it belongs to
\(\Sigma_R\).

Any \(A\)-\(A\) equality the first auxiliary key generates is
therefore a consequence of \(\Sigma_R\).

\smallskip\noindent\textbf{Step 7b --- Second auxiliary key, sub-case
\(m=1\).}
The determinant of the second key is the single position
\(P_\iota.2\).  Via the bidirectional inclusion, this is tied to
\(P_j.i_j\).  By the symmetric argument to Step~7a, the projection
gives the full key
\[
i_j\to P_j.1\,P_j.2
\]
on \(P_j\), also a member of \(\Sigma_R\).

\smallskip\noindent\textbf{Step 7c --- Second auxiliary key, sub-case
\(m>1\).}
The determinant spans positions \(\{2,\ldots,m+1\}\), which correspond
to \(m\) different identifying relationships \(P_1,\ldots,P_m\) (not
all equal to \(P_j=R\)).  The argument here is structural rather than
via direct projection.

By the origin invariant in
Lemma~\ref{lem:kger-key-nonkey-separation}, every \(P_\iota\)-fact has
its first component and at most one further component (some position
\(k+1\)) copied from an identifying \(P_k\)-fact; all remaining
non-first components are fresh witnesses introduced when that
\(P_\iota\)-fact was created.

Since the initial \(R\)-tableau contains only \(R\)-facts
(\(R=P_j\), not \(P_k\) for \(k\neq j\)), the positions \(k+1\) with
\(k\neq j\) of any \(P_\iota\)-fact are populated only by generated
fresh witnesses.  For the second key to fire on two \(P_\iota\)-facts,
they must agree on every position \(2,\ldots,m+1\), including those
\(k\neq j\) positions filled by fresh witnesses.

Such agreement on fresh-witness positions is mediated by generated
facts and prior key equalities outside the initial \(R\)-tableau.  By
Lemma~\ref{lem:kger-old-domain-equality}, this generated material
cannot create a new equality between distinct \(A\)-elements beyond
what \(\mathcal B_{\mathcal S}\cup\mathcal D_{\mathcal S}\cup
\mathcal K_{\mathcal S}\) alone forces.  Hence an application of the
second auxiliary key with \(m>1\) either:
\begin{itemize}
\item has no effect on \(A\)-variables at all, or
\item identifies first components that are already equal in the chase
quotient --- i.e., the equality it generates is already a consequence
of prior steps,
\end{itemize}
and in neither situation does it produce a new key consequence over
\(R\) beyond \(\Sigma_R\).

\smallskip\noindent\textbf{Step 8 --- Case C: key step on a
non-auxiliary \(Q\neq R\).}
The initial \(R\)-tableau contains no \(Q\)-facts (only \(R\)-facts
and unary type facts).  Such a key step could affect \(A\)-variables
only through \(Q\)-facts generated during the chase.

By Lemma~\ref{lem:kger-old-domain-equality}, adding
\(\mathcal N_{\mathcal S}\) before the key chase cannot create a new
\(A\)-\(A\) equality beyond those already forced by
\(\mathcal B_{\mathcal S}\cup\mathcal D_{\mathcal S}\cup
\mathcal K_{\mathcal S}\).  By Lemma~\ref{lem:kger-key-nonkey-separation},
every \(Q\)-fact in the chase has at least one fresh component (since
\(Q\neq R\) and the initial instance has no \(Q\)-facts).  So a key
step on \(Q\) compares two \(Q\)-facts each with a fresh side; each
generated equality has at least one fresh side and cannot equate two
distinct \(A\)-elements.

Hence \(Q\)-key steps contribute no new key consequence over \(R\).

\smallskip\noindent\textbf{Step 9 --- Conclude.}
By Cases A, B (sub-cases 7a, 7b, 7c), and C, every \(A\)-\(A\)
equality generated by a key step in the chase is either a
\(\Sigma_R\)-consequence or has no effect on \(A\)-variables.  The
restricted quotient on \(A\) --- which is what determines whether the
chase forces the distinguished tuple --- is therefore generated
entirely by \(\Sigma_R\)-equalities.

Chasing the JD tableau with \(\Sigma_R\) alone produces the same
restricted quotient on \(A\).  If the original chase forces the
distinguished tuple, so does the \(\Sigma_R\)-chase.  By soundness and
completeness of the chase for JD implication,
\[
\mathcal S\models J \quad\Longrightarrow\quad \Sigma_R\models J.
\]

Combined with the easy direction (Step~2), this gives the iff in the
theorem.  The corollary ``\(R\) is in 5NF'' is exactly the iff applied
as the textbook PJ/NF criterion.
\end{proof}

\noindent\textbf{Statement of Corollary~\ref{cor:kger-no-symmetry} (\colorbox{yellow}{KG-ER$_0$ cannot express symmetric binary relations}).}
Let \(\mathcal S\) be an ER-satisfiable KG-ER\(_0\) schema, let
\(R\) be a declared relationship type of arity \(2\), and assume that
\(\mathcal S\) admits a model \(M\) containing an \(R\)-tuple
\((\alpha,\beta)\) with \(\alpha\neq\beta\); equivalently,
\(\mathcal S\not\models\forall x,y.\ R(x,y)\to x=y\).
Then \(\mathcal S\) does not entail that \(R\) is symmetric:
\[
\mathcal S\not\models \forall x,y.\ R(x, y) \rightarrow R(y, x).
\]

\begin{proof}[Proof of Corollary~\ref{cor:kger-no-symmetry} (KG-ER$_0$ cannot express symmetric binary relations)]

\smallskip\noindent\textbf{Step 1 --- Reduce to the compiled schema.}
By Theorem~\ref{thm:kger-definitional-extension} (KG-ER\(_0^+\) as a
conservative extension), entailments over the original vocabulary are
preserved and reflected by passing to the compiled expansion
\(\mathcal S^+\).  Since symmetry of the original relationship \(R\)
is an original-vocabulary property (it only mentions \(R\), not any
auxiliary \(P_\iota\)), we may work in \(\mathcal S^+\) without loss.
This grants us the explicit auxiliary relationships and ordinary
key/inclusion machinery used throughout the development.

\smallskip\noindent\textbf{Step 2 --- Remark: necessity of the
non-diagonal hypothesis.}
The non-diagonal hypothesis is genuinely needed: every diagonal binary
relation satisfies \(R(x,y)\to R(y,x)\) vacuously (because \(x=y\)).
If \(\mathcal S\) forced \(R\) diagonal, then \(\mathcal S\) would
entail symmetry, and the corollary's conclusion would fail.  The
hypothesis rules out exactly this scenario.

\smallskip\noindent\textbf{Step 3 --- Proof by contradiction: chase
test setup.}
Assume for contradiction that
\(\mathcal S\models\forall x,y.\ R(x,y)\to R(y,x)\).  Use the standard
chase test for entailment of an inclusion-style sentence: take the
\emph{initial instance} \(I\) consisting of
\begin{itemize}
\item the single fact \(R(a,b)\), where \(a\) and \(b\) are fresh
distinct labelled nulls, and
\item the type facts forced by the declaration of \(R\) (entity/value
types on positions 1 and 2).
\end{itemize}
Set \(A=\{a,b\}\).  The chase test says: \(\mathcal S\) entails
\(R(x,y)\to R(y,x)\) if and only if, on every non-failing branch of
the chase of \(I\) with \(\mathcal S\), the quotient (after the key
chase) contains \(R(b,a)\).

If we can exhibit a single non-failing branch whose quotient does not
contain \(R(b,a)\), then symmetry is not entailed --- contradicting
the assumption.

\smallskip\noindent\textbf{Step 4 --- Building a non-failing branch
via the model \(M\).}
We use the model \(M\) from the hypothesis to choose a non-failing
branch.

Define \(h(a)=\alpha\) and \(h(b)=\beta\).  This extends to a
homomorphism \(h:I\to M\):
\begin{itemize}
\item \(R(a,b)\) maps to \(R(\alpha,\beta)\), which is in \(R^M\) by
hypothesis.
\item The type facts of \(I\) are forced by the declaration of \(R\)
(every model of \(\mathcal S\) satisfies them), so they map to type
facts true in \(M\).
\end{itemize}
So we have a homomorphism from \(I\) into a model of \(\mathcal S\).
Choose cover branches in the chase according to \(M\).  Along this
branch: existential rules pick witnesses consistent with \(M\);
disjointness constraints do not fire (they hold in \(M\)); the
subsequent key chase does not fail (keys hold in \(M\)).

This gives us a concrete non-failing branch to analyse.

\smallskip\noindent\textbf{Step 5 --- What must happen for symmetry to
hold.}
If \(\mathcal S\) entails symmetry, the quotient of this branch must
contain \(R(b,a)\): more precisely, some \(R\)-fact \(R(s,t)\) in the
chase, after equivalence closure by the key chase, must satisfy
\(s\equiv b\) and \(t\equiv a\).  In the language of
Lemma~\ref{lem:kger-old-domain-binary-fact}, the quotient must contain
an \(R\)-fact with \emph{ordered old-domain trace \((b,a)\)}.

We will show this is impossible along the chosen non-failing branch.
The argument has two stages, mirroring the architecture of the
lemmas:
\begin{itemize}
\item \emph{First stage:} the restricted chase (without
\(\mathcal N_{\mathcal S}\)) does not produce ordered trace \((b,a)\),
nor does it identify \(a\) with \(b\).
\item \emph{Second stage:} by
Lemma~\ref{lem:kger-old-domain-binary-fact}, the full chase (with
\(\mathcal N_{\mathcal S}\) added) inherits the same absence ---
adding \(\mathcal N_{\mathcal S}\) creates no new ordered
\(A^2\)-traces.
\end{itemize}

\smallskip\noindent\textbf{Step 6 ---
Lemma~\ref{lem:kger-key-nonkey-separation} restricts the all-old
\(R\)-facts.}
By Lemma~\ref{lem:kger-key-nonkey-separation} (chase-level key/non-key
separation), the non-key stage of the chase cannot create a new
\(R\)-fact whose components both lie in \(A=\{a,b\}\).  The only
\(R\)-fact in the chase that is entirely over \(A\) (before any key
step) is the original \(R(a,b)\) from \(I\).

So any candidate \(R\)-fact with ordered trace \((b,a)\) in the final
quotient must either:
\begin{itemize}
\item be the original \(R(a,b)\), after a key step has identified
\(a\) with \(b\) (in which case \(R(a,b)=R(b,a)\) trivially), or
\item be a generated \(R\)-fact \(R(s,t)\) with at least one fresh
component, brought into trace \((b,a)\) via key-step quotienting.
\end{itemize}

\smallskip\noindent\textbf{Step 7 --- Restricted chase
(\(\mathcal B_{\mathcal S}\cup\mathcal D_{\mathcal S}\cup
\mathcal K_{\mathcal S}\) only) does not produce \(R(b,a)\).}
Consider first the chase using only
\(\mathcal B_{\mathcal S}\cup\mathcal D_{\mathcal S}\cup
\mathcal K_{\mathcal S}\) (omitting \(\mathcal N_{\mathcal S}\)).

The \(\mathcal D_{\mathcal S}\)-inclusions can copy the binary
projection of \(R(a,b)\) into auxiliary relations \(P_\iota\) (via
forward Id-inclusions), and project back from \(P_\iota\) to \(R\)
(via reverse Id-inclusions).  Crucially, these projections
\emph{preserve the orientation}: a forward inclusion
\(R[\ols{i_j},i_j]\subseteq P_\iota[1,j+1]\) followed by a reverse
inclusion \(P_\iota[1,j+1]\subseteq R[\ols{i_j},i_j]\) reproduces
\(R(a,b)\) in its original orientation, not \(R(b,a)\).  Position
\(\ols{i_j}\) of the source \(R\)-fact ends up at \(P_\iota.1\), which
projects back to position \(\ols{i_j}\) of \(R\) --- same place.

Every other relationship fact generated by these inclusions contains
a fresh witness (by Lemma~\ref{lem:kger-key-nonkey-separation}'s
auxiliary origin invariant).

Now consider key steps in this restricted chase.  The only ways to
potentially equate \(a\) with \(b\) or to produce an \(R\)-fact in
orientation \((b,a)\) would be:
\begin{itemize}
\item a key on \(R\) comparing \(R(a,b)\) with itself (trivial) or
with a fresh-component \(R\)-fact (which can only merge a fresh
witness with \(a\) or \(b\), not \(a\) with \(b\));
\item a key on \(P_\iota\) propagating through the bidirectional
inclusion, but the projections preserve orientation;
\item a key on some other relationship, all of whose facts have fresh
components.
\end{itemize}
The restricted \(D\)-supported key chase therefore neither identifies
\(a\) with \(b\) nor produces an all-old \(R\)-fact with ordered
trace \((b,a)\).

\smallskip\noindent\textbf{Step 8 ---
Lemma~\ref{lem:kger-old-domain-binary-fact} extends the conclusion to
the full chase.}
Now add the remaining non-key constraints \(\mathcal N_{\mathcal S}\)
back into the picture.  By
Lemma~\ref{lem:kger-old-domain-binary-fact} (old-domain binary-fact
preservation), if the full chase quotient contained an all-old
\(R\)-fact with ordered trace \((b,a)\), then the restricted
(\(\mathcal B_{\mathcal S}\cup\mathcal D_{\mathcal S}\cup
\mathcal K_{\mathcal S}\)) chase quotient would also contain such an
\(R\)-fact.

By Step~7, the restricted chase quotient does not contain an
\(R\)-fact with ordered trace \((b,a)\).  Therefore, by
Lemma~\ref{lem:kger-old-domain-binary-fact}, neither does the full
chase quotient.

In particular, the final quotient omits \(R(b,a)\).

\smallskip\noindent\textbf{Step 9 --- Contradiction.}
The chosen non-failing branch of the chase has a final quotient that
omits \(R(b,a)\) --- contradicting the assumption that
\(\mathcal S\) entails symmetry of \(R\) (which would require
\(R(b,a)\) in the quotient of every non-failing branch).

Therefore \(\mathcal S\) cannot entail that \(R\) is symmetric.
\end{proof}

\newpage
\section*{Appendix A.2. Proofs for Section 2}

\noindent\textbf{Statement of Theorem~\ref{lem:carm-5}.}
Let \(\mathcal C\) be a satisfiable \textsf{CARM} schema, and let \(R\)
be a relation symbol that is non-empty in some model of \(\mathcal C\).
Let \(\Sigma_R\) be the set of all key dependencies
\(X\to\operatorname{att}(R)\) entailed by \(\mathcal C\). Then, for every
join dependency \(J\) over \(\operatorname{att}(R)\),
\[
\mathcal C\models J
\quad\text{iff}\quad
\Sigma_R\models J .
\]
Hence every such relation \(R\) is in 5$^{th}$ normal form.

\begin{proof}[Proof of Theorem~\ref{lem:carm-5}]
The right-to-left direction follows by transitivity of implication,
because every dependency in \(\Sigma_R\) is entailed by
\(\mathcal C\).

For the converse, use the standard chase test for implication of a join
dependency \(J\) over \(R\).  The chase starts from the usual JD tableau,
whose rows are projections of one distinguished \(R\)-tuple.  The chase
succeeds exactly when the distinguished tuple is forced by the
constraints.  Let \(A\) be the active domain of this initial tableau.
Since \(R\) is non-empty in some model of \(\mathcal C\), the tableau has
a homomorphism into a model of \(\mathcal C\) sending every row to the
same \(R\)-tuple.  Hence there is a non-failing chase branch, with cover
choices made according to that model.

We first isolate what the non-key constraints can do on such a branch.
An inclusion dependency copies one existing projection tuple into one
target projection occurrence and fills the remaining attributes of the
target relation, if any, with fresh witnesses.  A cover constraint is the
same operation with a branch choice from the target projection to one of
the source projections.  Disjointness constraints only close branches.
Thus every non-key step is controlled by a single projection tuple; it
never combines two projection tuples.

By induction on the non-key chase, every generated fact has a single
origin row of the initial \(R\)-tableau.  More precisely, its old-domain
components are copied from projection occurrences of one initial row, and
all remaining components are fresh witnesses.  The occurrence
non-overlap condition ensures that, within a fixed relation, old
components copied from different projection occurrences cannot be pasted
together into a larger all-old tuple.  The fixed ordering of projection
occurrences also prevents a copied tuple from being permuted.  Therefore
the non-key part cannot create the distinguished \(R\)-tuple over \(A\)
unless that tuple was already present in the initial tableau.

Suppose now that the full chase with \(\mathcal C\) succeeds.  Since the
non-key stage cannot by itself create the distinguished tuple, success is
due to equality-generating key steps.  Consider the equivalence relation
that these key steps induce on the variables of the initial
\(R\)-tableau.  A key can affect such variables only through generated
facts whose old-domain part is a projection occurrence.  By
key/occurrence compatibility, that occurrence is, relative to the key of
its relation, either equal to the key, properly contained in the key, or
disjoint from it.  In the disjoint case the key determinant contains no
old component from the occurrence.  In the proper-containment case the
determinant contains at least one additional component that was fresh
when the fact was generated.  Before the first merge of two distinct
old-domain classes, such a fresh component has either no old trace or a
single old trace already copied into it by earlier key steps.  Agreement
on the larger determinant therefore cannot be the first cause of a merge
between two distinct old-domain classes.  The only case that can force a
new equality among initial \(R\)-variables is the case in which the
copied occurrence is itself a key determinant.

Transported back along the single-origin chain to the initial
\(R\)-tableau, that last case is exactly a key implication over \(R\):
the standard two-row chase test with the corresponding determinant on
\(R\) forces equality of the whole two tuples.  Hence the induced
equality on the initial \(R\)-tableau is generated by key dependencies on
\(R\) entailed by \(\mathcal C\), all of which are included in
\(\Sigma_R\).  Chasing the JD tableau with \(\Sigma_R\) therefore
produces the same quotient on the initial variables and also forces the
distinguished tuple.  By soundness and completeness of the chase for
join-dependency implication, \(\Sigma_R\models J\).
\end{proof}

\newpage
\section*{Appendix A.3. Proofs for Section 3}

\noindent\textbf{Statement of Theorem~\ref{thm:kgerzero-to-carm-oids}.}
Let $\mathcal S$ be a KG-ER$_0$ schema.  Then
$\mathcal S$ can be translated in linear time into a \textsf{CARM}
schema $\mathcal C_{\mathcal S}$ over the same OID and value domains.
The models of $\mathcal S$ and the models of $\mathcal C_{\mathcal S}$
are in bijection, up to the definitional expansion/reduct of the
relations introduced for \textsc{Id} constraints.

\begin{proof}[Proof of Theorem~\ref{thm:kgerzero-to-carm-oids}]
\emph{Step 1: compile identities.}
For each identity constraint
\[
  \iota=\textsc{Id}(E\ (P_1.i_1,\ldots,P_m.i_m))   
\]
introduce a fresh relation symbol
$I_\iota(C,A_1,\ldots,A_m)$, where $C$ has type $E$ and
$A_j$ has type $\tau(P_j,i_j)$.  Thus the typing inclusions include
$I_\iota.C\subseteq E$ and
$I_\iota.A_j\subseteq\tau(P_j,i_j)$.  Replace the binary identifying
relationships by the definitional views
\[
  P'_j=\pi_{C,A_j} I_\iota \qquad (1\leq j\leq m),
\]
with the inverse renaming of the roles of $P_j$.  Add to $I_\iota$ the two
keys
\[
  C\to A_1\cdots A_m,
  \qquad
  A_1\cdots A_m\to C,
\]
and the inclusion $E\subseteq I_\iota.C$.  Any remaining occurrence of an
identifying $P_j$ is translated through the same definitional view.
This does not discard information from $P_j$: after the role renaming,
if $(c,a)\in P'_j$, then typing gives $c\in E$; the inclusion
$E\subseteq P.1$ gives a tuple of
$P=P'_1\bowtie\cdots\bowtie P'_m$ with first component $c$; and the key
$P.1\to P.2\cdots P.(m+1)$ forces its $(j+1)$-st component to be $a$.
Hence $P'_j=\pi_{C,A_j}I_\iota$ in every legal instance.

\emph{Step 2: translate the remaining constraints.}
Keep all non-identifying relationship predicates, entity types, and
value types.  A type declaration $R(T_1,\ldots,T_n)$ gives unary
inclusions $R.i\subseteq T_i$.  A mandatory participation constraint
becomes a unary inclusion.  A relationship key remains the corresponding
relational key.  If such a constraint mentions an identifying
relationship $P_j$, first replace $P_j$ by
$\pi_{C,A_j}I_\iota$; unary inclusions become unary inclusions over
$C$ or $A_j$, and keys of the binary view become keys of $I_\iota$
using $C\to A_1\cdots A_m$.  \textsc{IsA} becomes unary inclusion, and
\textsc{Cover} and \textsc{Disjoint} are copied over the same inclusion
sources.

\emph{Step 3: verify the \textsf{CARM} conditions.}
All inclusion dependencies produced in Step~2 are unary, so their sources
and targets are either identical or disjoint from other such occurrences;
with respect to keys, a unary occurrence is either contained in a key or
disjoint from it.  The new identity relation $I_\iota$ has only keys.  
Moreover, the join dependency
\[
  I_\iota=\pi_{C,A_1}I_\iota\bowtie\cdots\bowtie\pi_{C,A_m}I_\iota
\]
is implied by the key $C\to A_1\cdots A_m$: tuples agreeing on $C$ must
agree on all attributes.  Hence the identity relation is in 5$^{th}$
normal form.  From the KG-ER$_0$ definitions, the remaining relation
schemas have no join dependencies other than those implied by their
declared keys.  Thus
$\mathcal C_{\mathcal S}$ is a \textsf{CARM} schema.

\emph{Step 4: prove losslessness.}
Let $M$ be a model of $\mathcal S$.  Define $\Phi(M)$ by keeping all
non-identifying predicates unchanged and by interpreting each new
$I_\iota$ as the join relation occurring in the semantics of the
corresponding \textsc{Id} constraint.  The two displayed keys and the inclusion
$E\subseteq I_\iota.C$ hold by the \textsc{Id} constraint, so $\Phi(M)$ is a
model of $\mathcal C_{\mathcal S}$.

Conversely, let $N$ be a model of $\mathcal C_{\mathcal S}$.  Define
$\Psi(N)$ by keeping the old predicates and reconstructing every
identifying binary relationship as the corresponding projection of
$I_\iota$.  The key $C\to A_1\cdots A_m$ implies
$I_\iota=\bowtie_j\pi_{C,A_j}I_\iota$, and the key
$A_1\cdots A_m\to C$ gives the identifying direction of
\textsc{Id}.  Hence $\Psi(N)$ satisfies the original KG-ER$_0$
constraints.  By construction $\Psi(\Phi(M))=M$, and
$\Phi(\Psi(N))=N$ after expanding or reducing the definitional views.
No OID or value is changed, and the construction scans the schema once;
therefore the translation is linear and lossless.
\end{proof}

\noindent\textbf{Statement of Theorem~\ref{thm:kgerzero-to-carm-no-oids}.}
Let $\mathcal S$ be an identity resolving KG-ER$_0$ schema.  
Then $\mathcal S$ can be transformed into an OID-free \textsf{CARM}
schema $\mathcal C^\circ_{\mathcal S}$ over value domains only.  The
models of $\mathcal S$ and the models of
$\mathcal C^\circ_{\mathcal S}$ are in bijection up to renaming of OIDs
and definitional expansion/reduct of the fixed identifier relationships.
Once the identities $\iota_K$ have been fixed, the transformation is
linear in the size of the schema.

\begin{proof}[Proof of Theorem~\ref{thm:kgerzero-to-carm-no-oids}]
\emph{Step 1: choose canonical value codes for OIDs.}
Let \(H_{\mathcal S}\) be the Hasse diagram of entity types; its roots
are the \emph{kinds}.  For each
root $K$ of \(H_{\mathcal S}\), fix one value identity entailed by
$\mathcal S$:
\[
  \iota_K=\textsc{Id}
  (K\ (P^K_1.i^K_1,\ldots,P^K_{m_K}.i^K_{m_K})).
\]
For a root $K$, write
\[
  \mathbf A_K=(A^K_1,\ldots,A^K_{m_K})
  \qquad\text{with}\qquad
  \operatorname{dom}(A^K_j)=\tau(P^K_j,i^K_j).
\]
In every model $M$ of $\mathcal S$, the fixed identity $\iota_K$ defines
a map
\[
  \lambda^M_K(o)=(a_1,\ldots,a_{m_K})
\]
where $(o,a_j)$ belongs to the $j$-th identifying relationship, after
the harmless role renaming used in the semantics of \textsc{Id}.  The
dependency $K\subseteq P.1$ makes $\lambda^M_K$ total on $K^M$,
$P.1\to P.2\cdots P.(m_K+1)$ makes it functional, and
$P.2\cdots P.(m_K+1)\to P.1$ makes it injective.  Since different roots
are disjoint, the pair $(K,\lambda^M_K(o))$ is a canonical replacement
for the OID $o$.

\emph{Step 2: replace each OID position by the code of its kind.}
If $E$ is below the root $K$, replace the unary entity relation $E$ by
an OID-free relation $E^\circ(\mathbf A_E)$, where $\mathbf A_E$ is a
fresh copy of $\mathbf A_K$, and declare $\mathbf A_E$ to be a key.  For
each relationship $R$, replace every entity-typed role $R.i$ whose type
lies below $K$ by a fresh copy $\mathbf A_{R,i}$ of $\mathbf A_K$; leave
value-typed roles unchanged.  The fixed root identifier relationships
need not be base relations: they are definitional projections of
$K^\circ$.  Any other constraint mentioning one of them is translated
through that projection.

\emph{Step 3: translate constraints.}
A key of $R$ is translated by replacing each entity role occurring in it
by the whole corresponding identifier block.  A typing inclusion
$R.i\subseteq E$ becomes
$\mathbf A_{R,i}\subseteq\mathbf A_E$ when $R.i$ is entity-typed, and
remains unary when $R.i$ is value-typed.  A mandatory participation
$E\subseteq R.i$ becomes
$\mathbf A_E\subseteq\mathbf A_{R,i}$.  If
$E_1,\ldots,E_n$ are below the same root, then \textsc{IsA},
\textsc{Cover}, and \textsc{Disjoint} constraints are copied over the
corresponding identifier blocks.  The remaining \textsc{Id} constraints
are translated by the same replacement; their two identifying
dependencies become keys of the translated identity relation.

\emph{Step 4: check that the result is \textsf{CARM}.}
Every source or target of an inclusion dependency is either a whole
identifier block or a single value attribute.  Identifier blocks are
freshly copied and ordered, so two such occurrences are either equal by
construction or disjoint.  With respect to keys, each occurrence is
either a key, a subset of a translated key, or disjoint from all keys.
Covering and disjointness constraints are still over sources of
inclusions with common targets.  Finally, the replacement of an OID
attribute by its whole identifier block is a bijective recoding of legal
tuples, with the block treated uniformly in every translated key and
dependency.  A join dependency of a translated relation therefore
corresponds, by collapsing each identifier block back to its OID, to a
join dependency of the original relation; since the original dependency
is key-implied, so is the translated one.  No declared dependency splits
an identifier block into proper subblocks.  Thus
$\mathcal C^\circ_{\mathcal S}$ is a \textsf{CARM} schema and has no
OID-valued attributes.

\emph{Step 5: prove losslessness.}
Given a model $M$ of $\mathcal S$, define $M^\circ$ by replacing every
OID $o$ of kind $K$ by $\lambda^M_K(o)$ in all entity and relationship
relations, and by interpreting the translated constraints as in
Steps~2--3.  Injectivity of $\lambda^M_K$ preserves keys, while the
definition of $M^\circ$ preserves inclusions, covers, and disjointness.

Conversely, given a model $N$ of $\mathcal C^\circ_{\mathcal S}$,
construct a KG-ER$_0$ model $N^{-}$ as follows.  For every root $K$ and
every tuple $\mathbf a\in (K^\circ)^N$, introduce the canonical OID
$(K,\mathbf a)$.  If $E$ lies below $K$, put $(K,\mathbf a)$ in
$E^{N^{-}}$ exactly when $\mathbf a\in (E^\circ)^N$.  Reconstruct every
relationship by replacing each identifier block $\mathbf a$ of kind $K$
by the OID $(K,\mathbf a)$ and by leaving all value components
unchanged.  For each fixed root identity $\iota_K$, reconstruct the
identifier relationships by projecting $\mathbf a$ to its components.
The translated constraints ensure that all reconstructed tuples are
well typed and satisfy the original keys, mandatory participations,
identity constraints, inclusions, covers, and disjointness constraints.

For every $M$, the instance $(M^\circ)^{-}$ is isomorphic to $M$ by the
OID renaming $o\mapsto(K,\lambda^M_K(o))$.  For every $N$,
$(N^{-})^\circ=N$ by construction.  Hence the transformation is
lossless, up to OID renaming and the stated definitional views.
\end{proof}

\noindent\textbf{Statement of Theorem~\ref{thm:carm-to-kgerzero}.}
Let $\mathcal C$ be a \textsf{CARM} schema without OID domains.  Assume
that the inclusion, covering, and disjoint constraints of $\mathcal C$
are given by their forest presentation: the nodes are ordered
source/target projection occurrences and the edges are the inclusion
dependencies between them.  Then $\mathcal C$ can be transformed into a
KG-ER$_0$ schema $\mathcal K_{\mathcal C}$ such that the instances of
$\mathcal C$ are in bijection, up to renaming of the introduced OIDs,
with the instances of $\mathcal K_{\mathcal C}$.  In particular, the
transformation is lossless.

\begin{proof}[Proof of Theorem~\ref{thm:carm-to-kgerzero}]
Write a node of the forest presentation as
$\beta=(R;A_1,\ldots,A_m)$, where $R$ is a relation symbol of
$\mathcal C$ and $A_1,\ldots,A_m$ is the ordered source or target
projection represented by that node.  Equal occurrences are identified.
The \textsf{CARM} non-overlap condition ensures that, for each fixed
relation $R$, two such occurrences are either equal or disjoint, so they
can be replaced independently.

\emph{Schema construction.}
For every node $\beta$ introduce an entity type $E_\beta$.  If
$\gamma$ is the parent of $\beta$ in the forest, i.e.,
$\beta\subseteq\gamma$ is an inclusion dependency of $\mathcal C$, add
\[
  \textsc{IsA}((E_\beta)\ E_\gamma).
\]
Covering and disjointness constraints over sources
$\beta_1,\ldots,\beta_n$ with common target $\gamma$ become
\[
  \textsc{Cover}((E_{\beta_1},\ldots,E_{\beta_n})\ E_\gamma),
  \qquad
  \textsc{Disjoint}((E_{\beta_1},\ldots,E_{\beta_n})\ E_\gamma).
\]

For each root $B=(R;A_1,\ldots,A_m)$ of the forest, introduce binary
relationship types
\[
  I^B_j(E_B,V^B_j) \qquad (1\leq j\leq m),
\]
where $V^B_j$ is the value type corresponding to the domain of
attribute $A_j$.  Add the identity constraint
\[
  \textsc{Id}(E_B\ (I^B_1.2,\ldots,I^B_m.2)).
\]
All entity types below $B$ use this root identity, since they are
subtypes of $E_B$.

For each relation $R$ of $\mathcal C$, form a relationship type
$R^\dagger$ as follows.  Each forest node
$\beta=(R;A_1,\ldots,A_m)$ is replaced by one entity-typed role
$e_\beta$ of type $E_\beta$; attributes of $R$ not belonging to any
forest node are kept as value-typed roles.  Add the mandatory
participation constraint
\[
  \textsc{Mand}(E_\beta,R^\dagger.e_\beta)
\]
for every such occurrence $\beta$.  Hence $E_\beta$ is exactly the
projection of $R^\dagger$ on the role $e_\beta$.

Keys are translated structurally.  If $K$ is a key of $R$, replace in
$K$ every whole occurrence $\beta=(R;A_1,\ldots,A_m)$ contained in $K$
by the role $e_\beta$, and keep all remaining value attributes.  The
\textsf{CARM} key-overlap condition ensures that no key cuts through a
forest occurrence.  Add the resulting \textsc{Key} constraint on
$R^\dagger$.

\emph{From \textsf{CARM} instances to KG-ER$_0$ instances.}
Let $M$ be a model of $\mathcal C$.  For every root $B$ and every tuple
\[
  \mathbf a=(a_1,\ldots,a_m)
\]
occurring in the projection represented by $B$, introduce the canonical
OID $(B,\mathbf a)$.  For a node $\beta$ below $B$, put
$(B,\mathbf a)$ in $E_\beta$ exactly when $\mathbf a$ occurs in the
projection represented by $\beta$.  This satisfies the translated
\textsc{IsA}, \textsc{Cover}, and \textsc{Disjoint} constraints because
the corresponding projection inclusions, covers, and disjointness
constraints hold in $M$.

Populate the root identifier relationships by
$((B,\mathbf a),a_j)\in I^B_j$ for every root tuple $\mathbf a$ and
every $j$.  Finally, for each tuple $t\in R^M$, put in
$(R^\dagger)^{\Phi(M)}$ the tuple obtained by replacing each occurrence
block $t[A_1,\ldots,A_m]$ belonging to a node $\beta$ below root $B$ by
the OID $(B,t[A_1,\ldots,A_m])$, and by leaving all other value
components unchanged.  Mandatory participation holds by construction.
The translated keys hold because equality of canonical OIDs is exactly
equality of the corresponding source or target tuples, and the original
keys hold in $M$.

\emph{From KG-ER$_0$ instances to \textsf{CARM} instances.}
Let $N$ be a model of the constructed KG-ER$_0$ schema.  For an OID
$o\in E_B^N$ at a root $B$, the identity constraint for $E_B$ gives a
unique value tuple
\[
  \lambda_B^N(o)=(a_1,\ldots,a_m)
\]
with $(o,a_j)\in (I^B_j)^N$ for all $j$; the reverse key of the identity
constraint makes this tuple injective in $o$.  For a node $\beta$ below
$B$, use the same tuple $\lambda_B^N(o)$ for every $o\in E_\beta^N$.

Define a \textsf{CARM} instance $\Psi(N)$ by replacing, in every tuple
of every $(R^\dagger)^N$, each entity role $e_\beta$ containing an OID
$o$ below root $B$ by the tuple $\lambda_B^N(o)$, and by leaving all
value roles unchanged.

The instance $\Psi(N)$ satisfies the constraints of $\mathcal C$.  For
an inclusion $\beta\subseteq\gamma$, if $\mathbf a$ occurs in the
$\beta$-projection of $\Psi(N)$, then some OID $o$ with
$\lambda_B^N(o)=\mathbf a$ occurs in role $e_\beta$; hence
$o\in E_\beta^N$.  The translated \textsc{IsA} constraint gives
$o\in E_\gamma^N$, and mandatory participation of $E_\gamma$ gives a
tuple of the target relationship containing the same OID $o$, whose
identifier tuple is again $\mathbf a$.  Thus $\mathbf a$ occurs in the
target projection.  The same argument with translated \textsc{Cover}
proves covering.  For disjointness, if the same tuple occurred in two
disjoint source projections, injectivity of the root identity would make
the two witnessing OIDs equal, contradicting the translated
\textsc{Disjoint} constraint.

Keys are preserved as well.  If two tuples of a reconstructed relation
agree on a key of $\mathcal C$, then their corresponding KG tuples agree
on the translated key: for an entity role this follows from injectivity
of the root identity, which turns equality of identifier tuples into
equality of OIDs.  The translated \textsc{Key} constraint on
$R^\dagger$ therefore makes the KG tuples equal, and the reconstructed
CARM tuples are equal.

\emph{Comparison of the two maps.}
For every CARM model $M$, $\Psi(\Phi(M))=M$ by construction.  Conversely,
for every KG-ER$_0$ model $N$, $\Phi(\Psi(N))$ is isomorphic to $N$ via
the OID map
\[
  o\mapsto (B,\lambda_B^N(o))
\]
for OIDs below root $B$.  This map is injective by the root identity and
surjective because mandatory participation makes every entity OID occur
in its representative relationship role.  It preserves all translated
relationship tuples and all root identifier relationships.  Hence the
two schemas have the same models up to renaming of introduced OIDs, and
the transformation is lossless.
\end{proof}

\newpage
\section*{Appendix A.4. Proofs for Section 4}

\noindent\textbf{Statement of Theorem~\ref{thm:gcarm-to-carm}.}
Let $\mathcal G$ be a \textsf{gCARM} schema. By
Theorem~\ref{thm:fdmvd}, the lossless and dependency-preserving part of
such a decomposition exists whenever the FD/MVD constraints of $R$ have
an extended conflict-free cover.  Then $\mathcal G$ can be transformed
into a \textsf{CARM} schema $\mathcal C_{\mathcal G}$ whose instances
are in bijection with the instances of $\mathcal G$.

\begin{proof}[Proof of Theorem~\ref{thm:gcarm-to-carm}]
%For each relation $R(U)$ of $\mathcal G$, replace $R$ by the component
%relations $R^1(U^1),\ldots,R^k(U^k)$ of the chosen decomposition
%$\Delta_R$.  The keys of each component are the keys representing the
%dependencies projected on that component.  Add the
%projection-preservation inclusions belonging to the decomposition.
%Finally, translate every original inclusion, covering, and disjointness
%constraint as follows: if an ordered occurrence $X$ of $R$ occurs in such
%a constraint, interpret it as the projection $X$ of the representative
%component $R^{\operatorname{rep}_R(X)}$.
%
%The resulting schema is a \textsf{CARM} schema.  Each component relation
%is in 5$^{th}$ normal form by the chosen decomposition, and its remaining
%dependencies are represented as keys.  The only non-key dependencies are
%the projection-preservation inclusions and the translations of the
%original inclusions.  By \textsf{CARM}-admissibility, their sources and
%targets satisfy the key-overlap, occurrence-overlap, and lexicographic
%ordering conditions required by \textsf{CARM}; covering and disjointness
%constraints are still over sources of inclusions with common targets.
%
%Define the instance map $\Phi$ from $\mathcal G$ to
%$\mathcal C_{\mathcal G}$ by projection:
%\[
%  (R^a)^{\Phi(I)}=\pi_{U^a}(R^I)
%  \qquad
%  \text{for every }R\text{ and every }R^a\in\Delta_R .
%\]
%Since each $\Delta_R$ is dependency-preserving, the component keys hold.
%The projection-preservation inclusions hold by construction, and the
%translated inclusion, covering, and disjointness constraints hold because
%they are evaluated on projections of the same original relations.
%
%Conversely, define $\Psi$ from $\mathcal C_{\mathcal G}$ to
%$\mathcal G$ by natural join:
%\[
%  R^{\Psi(J)}=(R^1)^J\bowtie\cdots\bowtie (R^k)^J .
%\]
%The lossless part of $\Delta_R$ gives
%$R^I=\bowtie_a\pi_{U^a}(R^I)$ for every legal $\mathcal G$-instance
%$I$, hence $\Psi(\Phi(I))=I$.  In the other direction,
%\textsf{CARM}-admissibility gives
%\[
%  \pi_{U^a}\bigl((R^1)^J\bowtie\cdots\bowtie (R^k)^J\bigr)=(R^a)^J
%  \qquad (1\leq a\leq k),
%\]
%so $\Phi(\Psi(J))=J$.  Moreover, every original inclusion, covering, and
%disjointness constraint holds in $\Psi(J)$ because its chosen
%representative projection in $\Phi(\Psi(J))$ is exactly the component
%projection already constrained in $J$.  Thus $\Phi$ and $\Psi$ are
%inverse bijections on legal instances, which proves losslessness.
\end{proof}

\noindent\textbf{Statement of Theorem~\ref{thm:carm-to-gcarm}.}
Every \textsf{CARM} schema $\mathcal C$ can be viewed as a
\textsf{gCARM} schema $\mathcal G_{\mathcal C}$ over the same relation
symbols and domains.  The two schemas have exactly the same instances;
hence the transformation is lossless.

\begin{proof}[Proof of Theorem~\ref{thm:carm-to-gcarm}]
%Keep all relation symbols, keys, inclusion dependencies, covering
%constraints, and disjointness constraints of $\mathcal C$.  For each
%relation $R(U)$, let $D_R$ contain the FD $K\to U$ for every declared key
%$K$ of $R$.  In addition, for every ordered source or target occurrence
%$X$ of an inclusion, covering, or disjointness constraint on $R$, add the
%trivial FD $X\to X$.
%
%These added trivial dependencies do not change the legal instances.
%They only make the \textsf{gCARM} syntactic condition on inclusion
%occurrences explicit: every source or target occurrence is now equal to
%the source of some dependency.  The second non-overlap condition on
%different inclusion occurrences is already one of the \textsf{CARM}
%conditions, and the covering and disjointness constraints are copied
%unchanged.
%
%It remains to check the extended conflict-free requirement.  The
%non-trivial dependencies in each $D_R$ are keys.  If $X$ is a key source,
%then every MVD $X\twoheadrightarrow Y$ implied by $D_R$ is also an FD
%$X\to Y$, and so it is excluded from $\operatorname{Env}(D_R)$.  If $X$
%is the source of one of the added trivial FDs, then $D_R$ entails no
%genuine MVD with source $X$ beyond the trivial consequences of
%reflexivity.  Therefore $\operatorname{Env}(D_R)=\emptyset$, which is
%conflict-free.  Thus each $D_R$ has an extended conflict-free cover.
%
%The identity map on instances is therefore an isomorphism between
%$\mathcal C$ and $\mathcal G_{\mathcal C}$, so the transformation is
%lossless.
\end{proof}


\bibliographystyle{named}
\bibliography{KG-ER-zero.bib}
\end{document}
